\RequirePackage[svgnames,table]{xcolor}

\documentclass[11pt,letterpaper,logo]{yalearxiv}

\usepackage{graphicx}
\usepackage{float,epstopdf}
\usepackage{bbm}

\usepackage{microtype}

\usepackage{natbib}
\setcitestyle{square}

% \usepackage{subfig}
% \usepackage{subfigure}
\usepackage{subcaption}
\usepackage{booktabs}

\usepackage{amsmath}
\usepackage{amssymb}
\usepackage{mathtools}
\usepackage{amsthm}
\usepackage{dsfont}
\usepackage{multicol}
\usepackage{makecell}
\usepackage{multirow} 
\usepackage{amsfonts} 
\usepackage{mathrsfs}
\usepackage[amssymb, thickqspace]{SIunits}
\usepackage{enumitem}
\usepackage{pgfplotstable}
\usepackage{lipsum}		% Can be removed after putting your text content
% \usepackage[pagebackref,breaklinks,colorlinks]{hyperref}


\usepackage{microtype}
\usepackage{graphicx}
\usepackage{subcaption}
\usepackage{booktabs} % for professional tables
\usepackage[table]{xcolor}
\usepackage{arydshln}
% \newcolumntype{s}{>{\columncolor[HTML]{E9E9E9}}c}
\usepackage[normalem]{ulem} % [normalem] prevents the package from changing the default behavior of `\emph` to underline.



\usepackage{cases}
\usepackage{wrapfig}


% \usepackage{algorithm}
% \usepackage{algorithmic}
% \newcommand{\theHalgorithm}{\arabic{algorithm}}

% \renewcommand{\algorithmicrequire}{\textbf{Input:}} 
% \renewcommand{\algorithmicensure}{\textbf{Output:}}

\usepackage{url}

% Todonotes is useful during development; simply uncomment the next line
%    and comment out the line below the next line to turn off comments

% thmtools/thm-restate clash with the shared theorem counters in command.tex
% (\newtheorem{theorem}[definition]{...} -> "\c@theorem already defined").
% thmtools' aliasctr runs \@counteralias{theorem}{definition}, creating
% \c@theorem; thm-patch then hands the LaTeX 2026-06-01 kernel a *self*-alias
% \newcounteralias{theorem}{theorem}, whose \@ifdefinable guard rejects the
% now-existing \c@theorem. Overleaf pins an older kernel, so it builds there.
% A \RequirePackage[<older>]{latexrelease} rollback does NOT help. Neither
% package's features (\declaretheorem, restatable, \listoftheorems) is used
% anywhere in this paper, so the fix is simply to drop them.
% DO NOT let an Overleaf sync re-add these two lines.
% \usepackage{thmtools}
% \usepackage{thm-restate}
\usepackage{tabu}

\definecolor{huskypurple}{HTML}{4B2E83}
\newcommand{\ourclr}{purple}
% \usepackage[comma,authoryear,compress]{natbib}
\usepackage{titletoc}
\input{command}


\title{Kalman Delta Networks: Uncertainty-Aware Associative Memory}

\runningtitle{Kalman Delta Networks}

\usepackage{fontawesome5}

\definecolor{yaleblue}{RGB}{0,58,112}
\newcommand{\yale}{%
  \textsuperscript{%
    {\usefont{T1}{pbk}{m}{n}\textcolor{yaleblue}{\textbf{Y}}}%
  }%
}

\author{Ngoc Bui\yale, Rex Ying\yale\\
\yale Department of Computer Science, Yale University\\
\faGithub~\textbf{Source Code:} \href{https://github.com/ngocbh/kalman-linear-attention}{\texttt{https://github.com/ngocbh/kalman-linear-attention}}
}

\hypersetup{colorlinks=true, linkcolor=blue!50!black, citecolor=blue!50!black,
            urlcolor=blue!50!black}

% --- draft-local notation ---
\newcommand{\Fdelta}{\mathsf{F}}               % delta-rule residual block

% Abstract removed: the document class typesets \theabstract inside \maketitle,
% so define it empty (an empty abstract renders as nothing; no heading is added).
\providecommand{\theabstract}{}


\begin{document}


\maketitle


% =============================================================================
\section{Introduction}
\label{sec:intro}

\begin{itemize}[leftmargin=1.4em,itemsep=5pt,topsep=2pt]

  \item \textbf{Fixed-size memory, edited online.} Linear attention replaces the
  growing key--value cache of softmax attention with a fixed-size recurrent state
  $\ve{S}_t \in \R^{d_k \times d_v}$ that is \emph{written} as tokens stream in and
  \emph{read} by the query \citep{katharopoulos2020transformers},
  \begin{equation}
    \text{(write)}\quad \ve{S}_t = \ve{S}_{t-1} + \ve{k}_t \ve{v}_t^{\tr},
    \qquad\qquad
    \text{(read)}\quad \ve{o}_t = \ve{S}_t^{\tr} \ve{q}_t .
    \label{eq:linattn}
  \end{equation}
  The design question is how to \emph{edit} this store online --- what to write and
  what to forget. This plain additive write only ever \emph{adds} associations:
  nothing is removed and non-orthogonal keys interfere, so the store saturates and
  reads degrade.

  \item \textbf{DeltaNet: error-correcting writes.} The delta rule
  \citep{schlag2021linear,yang2024parallelizing} reads what is stored at $\ve{k}_t$ and
  writes back only the correction,
  \begin{equation}
  \begin{aligned}
    \ve{S}_t
      &= \ve{S}_{t-1} + \beta_t \ve{k}_t\,(\ve{v}_t-\ve{S}_{t-1}^{\tr}\ve{k}_t)^{\tr}
        \qquad\text{(canonical residual form)}\\
      &= (\ve{I}-\beta_t \ve{k}_t \ve{k}_t^{\tr})\,\ve{S}_{t-1}
         + \beta_t \ve{k}_t \ve{v}_t^{\tr},
  \end{aligned}
    \label{eq:deltanet}
  \end{equation}
  so a repeated key is \emph{overwritten} rather than accumulated. It still never
  forgets what \emph{other} keys hold.

  \item \textbf{Gated DeltaNet: forgetting by decay.} DeltaNet corrects the current key
  but removes nothing else, so stale associations linger and the state drifts over long
  contexts. To actually \emph{forget}, Gated DeltaNet shrinks the whole state by a
  data-dependent decay $\alpha_t\in(0,1)$ before each write \citep{yang2024gdn},
  \begin{equation}
  \begin{aligned}
    \ve{S}_t
      &= \alpha_t \ve{S}_{t-1}
         + \beta_t \ve{k}_t\,(\ve{v}_t-(\alpha_t \ve{S}_{t-1})^{\tr}\ve{k}_t)^{\tr}
        \qquad\text{(canonical residual form)}\\
      &= \alpha_t\,(\ve{I}-\beta_t \ve{k}_t \ve{k}_t^{\tr})\,\ve{S}_{t-1}
         + \beta_t \ve{k}_t \ve{v}_t^{\tr},
  \end{aligned}
    \label{eq:gdn}
  \end{equation}
  so untouched memories fade geometrically: the state stays bounded and recent tokens
  weigh more. But a single scalar $\alpha_t$ per head fades \emph{every} channel at the
  same rate --- it cannot hold some memories long while dropping others fast.

  \item \textbf{KDA: multi-timescale forgetting.} But a useful memory needs several
  horizons at once --- a name or topic kept for thousands of tokens, a local cue dropped
  immediately. KDA replaces the scalar with a diagonal $\ve{D}_t=\diag(\ve{\alpha}_t)$, giving each
  channel its own decay rate \citep{kimiteam2025kimilinear},
  \begin{equation}
  \begin{aligned}
    \ve{S}_t
      &= \ve{D}_t \ve{S}_{t-1}
         + \beta_t \ve{k}_t\,(\ve{v}_t-(\ve{D}_t \ve{S}_{t-1})^{\tr}\ve{k}_t)^{\tr}
        \qquad\text{(canonical residual form)}\\
      &= (\ve{I}-\beta_t \ve{k}_t \ve{k}_t^{\tr})\,\ve{D}_t\,\ve{S}_{t-1}
         + \beta_t \ve{k}_t \ve{v}_t^{\tr},
  \end{aligned}
    \label{eq:kda}
  \end{equation}
  so one head spans a \emph{mixture of timescales} --- the most expressive forgetting in
  the family, and the baseline we build on.

\end{itemize}


% =============================================================================
\section{Kalman Associative Memory}
\label{sec:memory-kalman}

\subsection{Associative Memory as a Linear-Gaussian State-Space Model}
\label{sec:memory-formulation}

The recurrent state in linear attention is not just a cache compression; it is an
\emph{estimate} of an associative memory. At time $t$, the memory is a matrix
$\ve{S}_t\in\R^{d_k\times d_v}$ that defines the associative map
\begin{equation}
  \mc M_t\colon \R^{d_k}\to\R^{d_v},\qquad \mc M_t(\ve{k})=\ve{S}_t^{\tr}\ve{k} .
  \label{eq:assoc-map}
\end{equation}
The query read is $\ve{o}_t=\mc M_t(\ve{q}_t)=\ve{S}_t^{\tr}\ve{q}_t$. The streaming token provides one
supervised measurement of this map: under key $\ve{k}_t$, the memory should return value $\ve{v}_t$,
i.e.\ $\mc M_t(\ve{k}_t)\approx \ve{v}_t$. This is exactly the shape of a state-space model: an
unobserved state evolves over time, and each token reveals a noisy linear projection of it.

\paragraph{Latent state.}
Let $\ve{S}_t^\star$ be the latent associative map the recurrent state is trying to track. The
transition describes how stored associations persist, decay, or drift before the next
measurement arrives:
\begin{equation}
  \ve{S}_t^\star = \ve{D}_t \ve{S}_{t-1}^\star + \ve{W}_t,\qquad \ve{W}_t\sim\N(0,\ve{\Omega}_t).
  \label{eq:mem-update-dynamics}
\end{equation}
The operator $\ve{D}_t$ is the process model: it is the formal object corresponding to forgetting
under information drift, topic drift, and other non-stationarity in the stream. For DeltaNet
$\ve{D}_t=\ve{I}$; for Gated DeltaNet $\ve{D}_t=\alpha_t \ve{I}$; for KDA
$\ve{D}_t=\diag(\ve{\alpha}_t)$. The process
noise $\ve{W}_t$ represents memory drift not explained by deterministic decay, while $\ve{\Omega}_t$ is
its covariance.

\paragraph{Observation.}
The current key reads one direction of the latent map. The observed value is
\begin{equation}
  \ve{v}_t = \ve{S}_t^{\star\tr}\ve{k}_t + \ve{e}_t,\qquad \ve{e}_t\sim\N(0,r_t \ve{I}).
  \label{eq:mem-update-observation}
\end{equation}
Thus $\ve{k}_t^{\tr}$ is the observation map and $\ve{v}_t$ is the measurement. The noise term $\ve{e}_t$
captures the part of the token value that should \emph{not} be treated as a clean memory
target. In an LLM, $\ve{v}_t$ is a contextual feature produced from a noisy token embedding,
local syntax, position, layer mixing, and finite model capacity; not all of it is a stable
fact that can be stored under key $\ve{k}_t$. The covariance $r_t\ve{I}$ sets how reliable this value
observation is: small $r_t$ means trust the token and write strongly, while large $r_t$ means
treat much of $\ve{v}_t$ as context-specific noise.

\paragraph{Filtering problem.}
The state-space model asks for an online filter: given the previous memory estimate and the
new noisy measurement $(\ve{k}_t,\ve{v}_t)$, update the associative memory so it tracks the
non-stationary stream. Forgetting is not an auxiliary gate added to attention; it is part of
the state dynamics that predict how memory changes before the next measurement arrives.


% =============================================================================
\subsection{Kalman Filtering}
\label{sec:kalman-filtering}

The state-space model in \S\ref{sec:memory-formulation} gives a principled solution to the
filtering problem: combine the prior memory estimate with each new measurement so that the
current state estimate has minimal uncertainty. Under the linear--Gaussian assumptions in
\eqref{eq:mem-update-dynamics} and \eqref{eq:mem-update-observation}, this becomes the
Kalman filtering update. The filter carries $\ve{S}_t$ as the current memory estimate, and a
covariance $\ve{P}_t\in\R^{d_k\times d_k}$ tracks uncertainty over key directions in that
estimate.

Tracking $\ve{P}_t$ is necessary because each token observes the associative map only at one key
direction $\ve{k}_t$, and the value $\ve{v}_t$ is noisy. After many tokens, some key directions are
well supported by past measurements, while others remain uncertain or have become stale
because of drift. High uncertainty means the current token should be allowed to edit the
memory strongly; low uncertainty means the prior memory should resist a noisy measurement.

\paragraph{Predict: forgetting before writing.}
Before observing $(\ve{k}_t,\ve{v}_t)$, the filter predicts how the state evolves from step $t-1$ to
step $t$:
\begin{equation}
  \widehat{\ve{S}}_t = \ve{D}_t \ve{S}_{t-1},
  \qquad
  \widehat{\ve{P}}_t = \ve{D}_t\ve{P}_{t-1}\ve{D}_t^{\tr}+\ve{\Omega}_t .
  \label{eq:kalman-predict}
\end{equation}
This is the formal role of forgetting. The transition $\ve{D}_t$ moves the associative memory
forward in time, while $\ve{\Omega}_t$ increases uncertainty for drift that the transition cannot
explain.

\paragraph{Read: compare the prediction to the token.}
The predicted value under the current key is
\begin{equation}
  \widehat{\ve{v}}_t = \widehat{\ve{S}}_t^{\tr}\ve{k}_t,
  \qquad
  \ve{\delta}_t = \ve{v}_t-\widehat{\ve{v}}_t .
  \label{eq:kalman-innovation}
\end{equation}
The vector $\ve{\delta}_t$ is the innovation: the part of the observed value that the predicted
memory does not explain.

\paragraph{Objective: minimize posterior uncertainty.}
After observing $(\ve{k}_t,\ve{v}_t)$, we update the predicted memory state by writing the innovation
back into memory. The only remaining choice is the write direction $\ve{\kappa}_t$: it controls
which key-space direction receives the residual. For a candidate $\ve{\kappa}_t$, the memory
update has the form
\begin{equation}
  \ve{S}_t=\widehat{\ve{S}}_t+\ve{\kappa}_t
  \big(\ve{v}_t-\widehat{\ve{S}}_t^{\tr}\ve{k}_t\big)^{\tr}.
  \label{eq:kalman-residual-write}
\end{equation}
The Kalman objective chooses the write direction that leaves the updated memory estimate as
certain as possible, i.e.\ the one that minimizes posterior uncertainty after the
measurement:
\begin{equation}
  \ve{\kappa}_t^\star
  =
  \arg\min_{\ve{\kappa}_t}
  \Trace\!\left(\ve{P}_t(\ve{\kappa}_t)\right),
  \qquad
  \ve{P}_t(\ve{\kappa}_t)
  =
  (\ve{I}-\ve{\kappa}_t \ve{k}_t^{\tr})\widehat{\ve{P}}_t
  (\ve{I}-\ve{\kappa}_t \ve{k}_t^{\tr})^{\tr}
  + r_t\,\ve{\kappa}_t\ve{\kappa}_t^{\tr}.
  \label{eq:kalman-tracking-objective}
\end{equation}
Here $\ve{P}_t(\ve{\kappa}_t)$ is the posterior covariance that would result if we used the
candidate write direction $\ve{\kappa}_t$. The trace is the total remaining uncertainty across
key-space directions, so minimizing it chooses the residual write that makes the filtered
memory estimate maximally certain after the token is used. The first term in
$\ve{P}_t(\ve{\kappa}_t)$ propagates the predicted uncertainty $\widehat{\ve{P}}_t$ through the correction
operator $\ve{I}-\ve{\kappa}_t\ve{k}_t^{\tr}$: observing key $\ve{k}_t$ should reduce uncertainty in the
directions corrected by the write. The second term, $r_t\ve{\kappa}_t\ve{\kappa}_t^{\tr}$, is the
uncertainty injected by trusting a noisy value measurement. Thus the objective balances two
effects: write strongly where the memory is uncertain, but avoid injecting too much
measurement noise when the token value is unreliable.

\paragraph{Update: write the uncertainty-weighted residual.}
Solving the objective in \eqref{eq:kalman-tracking-objective} gives the Kalman gain, which
decides how much to trust the token relative to the current memory:
\begin{equation}
  \ve{\kappa}_t = \frac{\widehat{\ve{P}}_t \ve{k}_t}
  {r_t+\ve{k}_t^{\tr}\widehat{\ve{P}}_t \ve{k}_t}.
  \label{eq:kalman-gain}
\end{equation}
The optimal filtering update writes only the innovation, along this uncertainty-weighted key
direction:
\begin{equation}
  \ve{S}_t = \widehat{\ve{S}}_t+\ve{\kappa}_t\ve{\delta}_t^{\tr},
  \qquad
  \ve{P}_t=(\ve{I}-\ve{\kappa}_t\ve{k}_t^{\tr})\widehat{\ve{P}}_t .
  \label{eq:kalman-update}
\end{equation}
Finally, the layer reads from the filtered memory,
\begin{equation}
  \ve{o}_t = \ve{S}_t^{\tr}\ve{q}_t .
  \label{eq:kalman-read}
\end{equation}
This filtering update decomposes memory editing into the three operations the model needs:
predict the drifting memory, measure the prediction error at the current key, and write back
the residual in proportion to uncertainty.

\refstepcounter{definition}\label{prop:kalman-optimal-update}
\begin{tcolorbox}[
  title={Proposition~\thedefinition: Kalman Optimal Update},
  colback=blue!7,
  colbacktitle=blue!22,
  colframe=blue!45!black,
  coltitle=black,
  fonttitle=\bfseries,
  boxrule=0.4pt,
  arc=1pt,
  left=5pt,
  right=5pt,
  top=5pt,
  bottom=5pt
]
Under the linear--Gaussian state-space model in
\eqref{eq:mem-update-dynamics}--\eqref{eq:mem-update-observation}, the residual write that
minimizes the posterior uncertainty in \eqref{eq:kalman-tracking-objective} is the Kalman
filtering update below. Given the memory estimate $\ve{S}_{t-1}$ and uncertainty $\ve{P}_{t-1}$
before token $t$, the filter first predicts how the memory drifts under the process model
$\ve{D}_t$ and process noise $\ve{\Omega}_t$. After observing the key--value pair
$(\ve{k}_t,\ve{v}_t)$, it writes only the value residual
$\ve{v}_t-\widehat{\ve{S}}_t^{\tr}\ve{k}_t$ along the uncertainty-weighted
direction $\ve{\kappa}_t$.
\begin{equation}
\begin{aligned}
  \text{predict:}\quad
  \widehat{\ve{S}}_t &= \ve{D}_t \ve{S}_{t-1},
  &
  \widehat{\ve{P}}_t &= \ve{D}_t \ve{P}_{t-1}\ve{D}_t^{\tr}+\ve{\Omega}_t,
  \\[3pt]
  \text{update:}\quad
  \ve{S}_t &= \widehat{\ve{S}}_t+\ve{\kappa}_t
  \big(\ve{v}_t-\widehat{\ve{S}}_t^{\tr}\ve{k}_t\big)^{\tr},
  &
  \ve{P}_t &= (\ve{I}-\ve{\kappa}_t \ve{k}_t^{\tr})\widehat{\ve{P}}_t .
  \\[3pt]
  \text{where}\quad
  \ve{\kappa}_t &= \frac{\widehat{\ve{P}}_t \ve{k}_t}
  {r_t+\ve{k}_t^{\tr}\widehat{\ve{P}}_t \ve{k}_t}.
\end{aligned}
\label{eq:kalman-optimal-update-box}
\end{equation}
\end{tcolorbox}

\noindent\textit{Proof.} See Appendix~\ref{app:proofs}.

\subsection{Delta-rule models as fixed-gain Kalman filters}
\label{sec:delta-family-kalman}

The update above separates two modeling choices: the transition $\ve{D}_t$, which predicts how
the memory changes before the token is written, and the gain $\ve{\kappa}_t$, which decides the
key-space direction and strength of the residual write. Delta-rule linear attention is the
fixed-gain special case of this filter. If we discard the covariance recursion and use an
isotropic surrogate $\widehat{\ve{P}}_t\approx b_t \ve{I}$ for the predicted uncertainty, then
\begin{equation}
  \ve{\kappa}_t
  =
  \frac{\widehat{\ve{P}}_t \ve{k}_t}{r_t+\ve{k}_t^{\tr}\widehat{\ve{P}}_t \ve{k}_t}
  \approx
  \frac{b_t}{r_t+b_t\|\ve{k}_t\|_2^2}\,\ve{k}_t
  \equiv
  \beta_t \ve{k}_t .
  \label{eq:fixed-gain-approx}
\end{equation}
For normalized keys, the exact Kalman gain therefore collapses to the scalar write strength
$\beta_t$ used by the delta-rule family. Substituting $\ve{\kappa}_t=\beta_t \ve{k}_t$ into the
Kalman update gives the shared fixed-gain residual form
\begin{equation}
\begin{aligned}
  \ve{S}_t
  &= \ve{D}_t \ve{S}_{t-1}
     + \beta_t \ve{k}_t
     \big(\ve{v}_t-(\ve{D}_t\ve{S}_{t-1})^{\tr}\ve{k}_t\big)^{\tr} \\
  &= (\ve{I}-\beta_t \ve{k}_t \ve{k}_t^{\tr})\ve{D}_t\ve{S}_{t-1}
     + \beta_t \ve{k}_t \ve{v}_t^{\tr}.
\end{aligned}
\label{eq:fixed-gain-delta-family}
\end{equation}
DeltaNet, Gated DeltaNet, and KDA are obtained by choosing different process models:
\begin{equation}
  \text{DeltaNet: }\ve{D}_t=\ve{I},\qquad
  \text{GDN: }\ve{D}_t=\alpha_t \ve{I},\qquad
  \text{KDA: }\ve{D}_t=\diag(\ve{\alpha}_t).
  \label{eq:delta-family-transitions}
\end{equation}
Thus the three models share the same measurement model and residual write; they differ only
in how the memory is predicted before the write. Relative to the Kalman optimal update, they
freeze the uncertainty dynamics and replace the adaptive gain with the fixed first-order
direction $\beta_t \ve{k}_t$. This does not require assuming $\ve{\Omega}_t=0$; rather, the delta-rule
approximation ignores the covariance prediction
$\widehat{\ve{P}}_t=\ve{D}_t\ve{P}_{t-1}\ve{D}_t^{\tr}+\ve{\Omega}_t$ altogether, so the gain no longer depends on the
process model $\ve{D}_t$ or the drift uncertainty $\ve{\Omega}_t$.


% =============================================================================
\section{Kalman Delta Networks}
\label{sec:kla}

The Kalman optimal update in \eqref{eq:kalman-optimal-update-box} is the target, but the
exact recursion is not directly compatible with fully parallel linear attention. Parallel
training relies on a scan form
\begin{equation}
  \ve{S}_t=\ve{A}_t\ve{S}_{t-1}+\ve{b}_t,
  \label{eq:parallel-affine-scan}
\end{equation}
where the per-token coefficients $\ve{A}_t$ and $\ve{b}_t$ are input-only, or are computable from
associative prefix statistics. Exact Kalman filtering breaks this structure because
$\ve{P}_t$ follows a Riccati recursion and the gain $\ve{\kappa}_t$ depends on the accumulated
posterior uncertainty. A Kalman Delta Network therefore asks which parts of the optimal
update can be relaxed while preserving the scan.

\paragraph{Diagonal information as the tractable state.}
We restrict the uncertainty state to the space of diagonal covariance matrices, the simplest
family that is still more expressive than a scalar gain:
\begin{equation}
  \ve{P}_t\in\{\diag(\ve{p}):\ve{p}\in\R_+^{d_k}\},
  \qquad
  \ve{P}_t=\diag(\ve{p}_t),
  \qquad
  \ve{C}_t=\ve{P}_t^{-1}=\diag(\ve{c}_t),
  \label{eq:diag-cov-info}
\end{equation}
where $\ve{p}_t$ is per-key-direction uncertainty and $\ve{c}_t$ is per-key-direction information.
This is the state-space analogue of diagonal RLS in Appendix~\ref{app:mem-rls}: instead of
learning one scalar write strength, the layer tracks which key channels are already well
measured and which remain uncertain.

\paragraph{Transition-aware diagonal prediction.}
With a diagonal process model $\ve{D}_t=\diag(\ve{\alpha}_t)$ and diagonal process noise
$\ve{\Omega}_t=\diag(\ve{\omega}_t)$, the Kalman covariance prediction remains pointwise:
\begin{equation}
  \widehat{\ve{p}}_t = \ve{\alpha}_t^2\odot \ve{p}_{t-1}+\ve{\omega}_t .
  \label{eq:diag-covariance-prediction}
\end{equation}
The gain is then the diagonal approximation of the Kalman gain,
\begin{equation}
  \ve{\kappa}_t =
  \frac{\widehat{\ve{p}}_t\odot \ve{k}_t}
  {r_t+\sum_i \widehat{\ve{p}}_{t,i}\ve{k}_{t,i}^2}.
  \label{eq:diag-kalman-gain}
\end{equation}
Unlike the fixed-gain delta rule, in the exact prediction~\eqref{eq:diag-covariance-prediction}
the gain $\ve{\kappa}_t$ depends on the same process model that predicts the memory: if a channel
is forgotten or assigned high process noise, its uncertainty changes before the write is formed.
How much of this coupling survives once the recursion is made scan-computable depends on the
parameterization, discussed below.

\paragraph{Parallel gain from M\"obius prefix statistics.}
A direct recurrent implementation would update $\ve{p}_t$ from the posterior covariance at
every token. The diagonal Kalman Delta Network instead keeps the information-form sufficient statistics
scan-computable, as in diagonal RLS, while allowing the process noise $\ve{\omega}_t$ to be
parameterized independently. In information variables,
\eqref{eq:diag-covariance-prediction} becomes
\begin{equation}
  \widehat{\ve{c}}_t
  =
  \frac{\ve{c}_{t-1}}
  {\ve{\alpha}_t^2+\ve{\omega}_t\odot\ve{c}_{t-1}},
  \qquad
  \widehat{\ve{p}}_t=\frac{1}{\widehat{\ve{c}}_t},
  \label{eq:prefix-diag-uncertainty}
\end{equation}
where products, ratios, and inverses are coordinatewise. This prediction step is exact under the
diagonal covariance restriction. The approximation enters at the measurement update. In the full
linear-Gaussian filter, observing value $\ve{v}_t$ through key $\ve{k}_t$ adds the rank-one Fisher
information of the observation,
\begin{equation}
  \ve{C}_t^{\mathrm{exact}}
  =
  \widehat{\ve{C}}_t+\frac{1}{r_t}\ve{k}_t\ve{k}_t^{\tr},
  \qquad
  \widehat{\ve{C}}_t=\diag(\widehat{\ve{c}}_t).
  \label{eq:diag-exact-info-update}
\end{equation}
Even when $\widehat{\ve{C}}_t$ is diagonal, the measurement term
$\ve{k}_t\ve{k}_t^{\tr}/r_t$ is dense, so an exact update would leave the diagonal family and
would no longer be represented by a vector $\ve{c}_t$. The diagonal KDN therefore projects the measurement
information back to a diagonal information state. The literal diagonal projection would add
$(\ve{k}_t\odot\ve{k}_t)/r_t$, but this is poorly scaled for $\ell_2$-normalized keys: a
spread-out coordinate has variance $1/d_k$, so the raw per-coordinate Fisher increment is
only $O(1/(d_k r_t))$. We instead use the standardized coordinate
$\widetilde{k}_{t,i}=\sqrt{d_k}\,k_{t,i}$ inside the diagonal projection, giving the
dimension-stable diagonal approximation
\begin{equation}
  \ve{u}_t=\frac{d_k}{r_t}\,\ve{k}_t\odot\ve{k}_t,
  \qquad
  \ve{c}_t=\widehat{\ve{c}}_t+\ve{u}_t,
  \qquad
  \ve{c}_0=\mu\ve{1}.
  \label{eq:diag-info-update}
\end{equation}
For $\ell_2$-normalized keys with spread-out coordinates, $\mathbb{E}[d_k k_{t,i}^2]\approx 1$,
so each channel receives an $O(1/r_t)$ information increment. The factor is $d_k$ because the
Fisher increment is quadratic in the scalar feature: standardizing the coordinate by
$\sqrt{d_k}$ multiplies its squared information by $d_k$. The scale conditions only the
uncertainty tracker used for future gains; the current residual-write gain still uses the
observation denominator in \eqref{eq:diag-kalman-gain}.
Combining prediction and update gives, for each channel $i$, a linear-fractional recurrence:
\begin{equation}
  c_{t,i}
  =
  \frac{c_{t-1,i}}
  {\alpha_{t,i}^2+\omega_{t,i}c_{t-1,i}}
  +u_{t,i}
  =
  \frac{(1+u_{t,i}\omega_{t,i})c_{t-1,i}
        +u_{t,i}\alpha_{t,i}^2}
       {\omega_{t,i}c_{t-1,i}+\alpha_{t,i}^2}.
  \label{eq:diag-info-mobius}
\end{equation}
Thus each token defines a M\"obius map per channel,
\begin{equation}
  c_{t,i}=f_{t,i}(c_{t-1,i}),
  \qquad
  f_{t,i}(c)=\frac{A_{t,i}c+B_{t,i}}{C_{t,i}c+D_{t,i}},
  \qquad
  \ve{M}_{t,i}
  =
  \begin{bmatrix}
    1+u_{t,i}\omega_{t,i} & u_{t,i}\alpha_{t,i}^2\\
    \omega_{t,i} & \alpha_{t,i}^2
  \end{bmatrix}.
  \label{eq:diag-info-mobius-matrix}
\end{equation}
M\"obius maps compose by matrix multiplication, so all prefix information states can be computed
by an associative scan over these $2\times2$ matrices:
\begin{equation}
  \begin{bmatrix}n_{t,i}\\ d_{t,i}\end{bmatrix}
  =
  \ve{M}_{t,i}
  \begin{bmatrix}n_{t-1,i}\\ d_{t-1,i}\end{bmatrix},
  \qquad
  c_{t,i}=\frac{n_{t,i}}{d_{t,i}},
  \qquad
  \begin{bmatrix}n_{0,i}\\ d_{0,i}\end{bmatrix}
  =
  \begin{bmatrix}\mu\\ 1\end{bmatrix}.
  \label{eq:diag-info-scan}
\end{equation}
This keeps the gain transition-aware without imposing the state-proportional restriction
$\omega_t\propto p_{t-1}$ required by the scalar affine scan.

\paragraph{Parameterization.}
The learnable per-token quantities should be the memory transition and the additive process
noise. Let $\ve{x}_t$ denote the current token representation at this layer. A simple
parameterization is
\begin{equation}
  \begin{aligned}
    \ve{\alpha}_t&=\sigma(\ve{W}_{\alpha}\ve{x}_t+\ve{b}_{\alpha}),
    &
    \ve{\omega}_t&=\omega_{\min}+\operatorname{softplus}(\ve{W}_{\omega}\ve{x}_t+\ve{b}_{\omega}),
    \\[2pt]
    r_t&=r_{\min}+\operatorname{softplus}(\ve{w}_{r}^{\tr}\ve{x}_t+b_r),
    &
    \ve{D}_t&=\diag(\ve{\alpha}_t).
  \end{aligned}
  \label{eq:kla-parameterization}
\end{equation}
The vector $\ve{\alpha}_t$ parameterizes the deterministic memory transition, while
$\ve{\omega}_t$ parameterizes absolute process drift in covariance units. In this sense,
$\ve{\omega}_t$ takes over the learnable role of ``how strongly should this token edit memory?''
from the scalar write gate $\beta_t$ used in GDN and KDA, but indirectly: it changes predicted
uncertainty, and then the gain $\ve{\kappa}_t$ becomes the actual adaptive write direction and
strength. The scalar $r_t$ is the token-dependent observation noise: it controls how reliable
the current value $\ve{v}_t$ is as a memory target. The floors $\omega_{\min}$, $r_{\min}$, and
$\mu$ can be fixed positive scalars or learned positive scalars per head; $\mu$ is used only to
initialize $\ve{c}_0$, not as a per-token information injection.

\paragraph{Affine memory scan.}
Once $\ve{\kappa}_t$ is computed from prefix statistics, the memory update is again an affine
scan:
\begin{equation}
\begin{aligned}
  \ve{S}_t
  &= \ve{D}_t\ve{S}_{t-1}
     + \ve{\kappa}_t\big(\ve{v}_t-(\ve{D}_t\ve{S}_{t-1})^{\tr}\ve{k}_t\big)^{\tr} \\
  &= (\ve{I}-\ve{\kappa}_t \ve{k}_t^{\tr})\ve{D}_t\ve{S}_{t-1}
     + \ve{\kappa}_t \ve{v}_t^{\tr}.
\end{aligned}
\label{eq:kla-diag-affine-scan}
\end{equation}

\refstepcounter{definition}\label{prop:kalman-linear-attention}
\begin{tcolorbox}[
  title={Proposition~\thedefinition: Diagonal Kalman Delta Network},
  colback=blue!7,
  colbacktitle=blue!22,
  colframe=blue!45!black,
  coltitle=black,
  fonttitle=\bfseries,
  boxrule=0.4pt,
  arc=1pt,
  left=5pt,
  right=5pt,
  top=5pt,
  bottom=5pt
]
The diagonal Kalman Delta Network is a scan-compatible diagonal-information approximation to the
Kalman optimal update. Given token representation $\ve{x}_t$, key--value pair
$(\ve{k}_t,\ve{v}_t)$, and previous memory $\ve{S}_{t-1}$, the layer parameterizes a
diagonal memory transition, additive process noise, and token-dependent observation noise,
\begin{equation}
\begin{aligned}
  \ve{\alpha}_t &= \sigma(\ve{W}_{\alpha}\ve{x}_t+\ve{b}_{\alpha}),
  &
  \ve{\omega}_t &= \omega_{\min}+\operatorname{softplus}(\ve{W}_{\omega}\ve{x}_t+\ve{b}_{\omega}),
  &
  r_t &= r_{\min}+\operatorname{softplus}(\ve{w}_{r}^{\tr}\ve{x}_t+b_r),
  \\[2pt]
  \ve{D}_t &= \diag(\ve{\alpha}_t).
\end{aligned}
\label{eq:kla-box-parameterization}
\end{equation}
It then computes transition-aware prefix information statistics by a per-channel M\"obius scan.
With $\ve{u}_t=d_k(\ve{k}_t\odot\ve{k}_t)/r_t$, define for each channel $i$
\begin{equation}
  \ve{M}_{t,i}
  =
  \begin{bmatrix}
    1+u_{t,i}\omega_{t,i} & u_{t,i}\alpha_{t,i}^2\\
    \omega_{t,i} & \alpha_{t,i}^2
  \end{bmatrix},
  \qquad
  \begin{bmatrix}n_{t,i}\\ d_{t,i}\end{bmatrix}
  =
  \ve{M}_{t,i}
  \begin{bmatrix}n_{t-1,i}\\ d_{t-1,i}\end{bmatrix},
  \qquad
  c_{t,i}=\frac{n_{t,i}}{d_{t,i}},
  \qquad
  \begin{bmatrix}n_{0,i}\\ d_{0,i}\end{bmatrix}
  =
  \begin{bmatrix}\mu\\ 1\end{bmatrix}.
  \label{eq:kla-box-mobius}
\end{equation}
The adaptive gain at token $t$ uses the exclusive prefix state $\ve{c}_{t-1}$:
\begin{equation}
\begin{aligned}
  \widehat{\ve{c}}_t
  &=
  \frac{\ve{c}_{t-1}}
  {\ve{\alpha}_t^2+\ve{\omega}_t\odot\ve{c}_{t-1}},
  &
  \widehat{\ve{p}}_t
  &= \frac{1}{\widehat{\ve{c}}_t},
  \\[3pt]
  \ve{\kappa}_t
  &=
  \frac{\widehat{\ve{p}}_t\odot\ve{k}_t}
  {r_t+\sum_i \widehat{\ve{p}}_{t,i}\ve{k}_{t,i}^2}.
\end{aligned}
\label{eq:kla-box-gain}
\end{equation}
With $\ve{\kappa}_t$ available before the memory scan, the memory update is an affine scan:
\begin{equation}
\begin{aligned}
  \ve{S}_t
  &=
  (\ve{I}-\ve{\kappa}_t\ve{k}_t^{\tr})\ve{D}_t\ve{S}_{t-1}
  +\ve{\kappa}_t\ve{v}_t^{\tr},
  &
  \ve{o}_t &= \ve{S}_t^{\tr}\ve{q}_t .
\end{aligned}
\label{eq:kla-box-memory}
\end{equation}
Thus the diagonal KDN preserves the Kalman structure of predicting memory through $\ve{D}_t$ and writing
an uncertainty-weighted residual, while replacing the dense covariance recursion by
parallel prefix statistics.
\end{tcolorbox}

Thus the diagonal approximation preserves the two ingredients missing from fixed-gain
delta-rule models while keeping the training primitive unchanged: the memory is predicted
through $\ve{D}_t$, and the residual write is weighted by a transition-aware uncertainty
estimate.

\paragraph{Reduction to fixed-gain delta rules.}
The gain \eqref{eq:diag-kalman-gain} collapses to a fixed-direction delta-rule write exactly when
the predicted uncertainty is isotropic, $\widehat{\ve{P}}_t\approx b_t\ve{I}$, equivalently
$\widehat{\ve{p}}_t=b_t\ve{1}$: then
$\ve{\kappa}_t=\frac{b_t}{r_t+b_t\|\ve{k}_t\|_2^2}\ve{k}_t\equiv\beta_t\ve{k}_t$ and
\eqref{eq:kla-diag-affine-scan} becomes
$\ve{S}_t=(\ve{I}-\beta_t\ve{k}_t\ve{k}_t^{\tr})\ve{D}_t\ve{S}_{t-1}+\beta_t\ve{k}_t\ve{v}_t^{\tr}$,
the vanilla delta-rule residual update with per-token write gate $\beta_t$ and transition
$\ve{D}_t=\diag(\ve{\alpha}_t)$. The transition alone then selects the baseline, with the
\emph{same} isotropic gain in every case: $\ve{\alpha}_t=\ve{1}$ gives DeltaNet, a uniform
$\ve{\alpha}_t$ gives GatedDeltaNet, and a \emph{per-channel} $\ve{\alpha}_t$ gives KDA. For KDA,
this means choosing $\ve{\omega}_t$ so that
$\widehat{\ve{p}}_{t,i}=\alpha_{t,i}^2 p_{t-1,i}+\omega_{t,i}=b_t$ for every channel, equivalently
$\omega_{t,i}=b_t-\alpha_{t,i}^2 p_{t-1,i}$ when this is nonnegative. The diagonal KDN strictly generalizes these
fixed-gain rules by allowing $\widehat{\ve{p}}_t$ to be anisotropic, so the write strength can
vary across key channels with the accumulated evidence.


% =============================================================================
\section{Experiments}
\label{sec:experiments}

\subsection{Synthetic Experiments}
\label{sec:synthetic-experiments}

We evaluate Kalman Delta Networks on the synthetic long-context tests used by
Kimi Linear~\citep{kimiteam2025kimilinear}. These tasks isolate the finite-memory
behaviors that motivate KDNs: exact copying, associative recall, and state tracking under
many interleaved updates. Following their setup, all synthetic runs use a small language
model with $2$ layers, $2$ attention heads, and head dimension $128$. For each task, we train
for at most $20{,}000$ steps and select the best training curve over the learning-rate grid
$\{5\times 10^{-5}, 10^{-4}, 5\times 10^{-4}, 10^{-3}\}$. We report two views of the same
protocol: peak accuracy as the sequence length increases from $256$ to $2048$ tokens, and
convergence speed at a fixed context length of $1024$ tokens.

\paragraph{Palindrome.}
The palindrome task tests exact copying from a compressed recurrent memory. The model is
given a random token sequence followed by a separator and must emit the sequence in reverse
order. For example, after observing
\[
  \texttt{O G R S U N E <sep>},
\]
the target continuation is
\[
  \texttt{E N U S R G O}.
\]
This task stresses whether the attention mechanism can preserve token-level information
across the whole prefix rather than only maintaining coarse statistics.

\paragraph{Multi-query associative recall.}
Multi-query associative recall (MQAR) tests whether the model can bind multiple keys to
their values and retrieve several queried values at the end of the context. A typical input
contains interleaved key--value pairs, then a separator, then query keys; for example,
\[
  \texttt{A 1 C 3 B 0 M 8 G 5 E 4 <sep> B G},
\]
with target outputs
\[
  \texttt{0 5}
\]
at the query positions. This task directly probes interference in fixed-size associative
memory: the model must retain the relevant bindings while ignoring unrelated pairs.

\paragraph{Stack state tracking.}
The state-tracking task follows the stack benchmark used in Kimi Linear. The stream contains
operations over $64$ independent last-in-first-out stacks, each identified by an ID. A push
operation such as \texttt{<push> 1 G} adds token \texttt{G} to stack $1$, while a pop query for
stack $0$ asks the model to predict the most recently pushed token on that stack. The model
must therefore maintain many evolving latent states and update only the addressed stack at
each operation. This setting tests whether a KDN can selectively overwrite stale information
without erasing unrelated memory channels.


\subsection{Pretraining Experiments}
\label{sec:pretraining-experiments}

\paragraph{Setup.}
Beyond the synthetic probes, we pretrain Kalman Delta Networks and competing linear-attention
mixers from scratch on FineWeb-Edu at two scales: a $220$M-parameter model trained on $20$B
tokens, used for a controlled mixer ablation, a $750$M-parameter model trained on $50$B
tokens, and a $1.3$B-parameter model trained on $100$B tokens. At each scale all mixers share the same backbone, optimizer, and data order, and are
matched both in parameter count---the feed-forward width is calibrated per mixer---and in
recurrent-state size, so that measured differences isolate the update rule rather than model
capacity. We train with AdamW (peak learning rate $4\times10^{-4}$, $\beta=(0.9,0.95)$, weight
decay $0.1$), gradient clipping at $1.0$, a cosine schedule with a $0.5$B-token warm-up, a global
batch of $0.5$M tokens, and a sequence length of $4$K. We compare against
DeltaNet~\citep{schlag2021linear}, Gated DeltaNet~\citep{yang2024gdn},
GDN-2~\citep{hatamizadeh2026gdn2}, and KDA~\citep{kimiteam2025kimilinear}; the Kalman Delta
Networks are the Isotropic and Diagonal variants developed above. We evaluate zero-shot with the
standard harness for recurrent language models and report in-context retrieval on RULER.

\paragraph{Language modeling and commonsense reasoning.}
Table~\ref{tab:pretrain-lm} reports WikiText and LAMBADA perplexity, zero-shot LAMBADA accuracy,
and the commonsense suite (PIQA, HellaSwag, WinoGrande, ARC-easy, ARC-challenge). Because the
recurrent-state size is matched, gains reflect a stronger update rule rather than a larger memory.
At $220$M/$20$B the two KDN variants and KDA form a clear top group, improving average accuracy
over GDN-2, Gated DeltaNet, and DeltaNet by roughly two points, with the Isotropic KDN attaining
the best perplexity and the best average. The ordering sharpens at $750$M/$50$B, where the Diagonal KDN
takes most columns --- LAMBADA perplexity and accuracy, WinoGrande, ARC-easy and ARC-challenge ---
while the Isotropic KDN takes WikiText and HellaSwag and the three Kalman variants tie on PIQA; the
Diagonal KDN holds the best average ($54.97$, versus $54.41$ for the Isotropic KDN, $53.87$ for KDA
and $51.45$ for GDN-2).
At $1.3$B/$100$B all three uncertainty-gated mixers beat GDN-2 on \emph{every} column --- average
accuracy rises from $58.51$ to $60.06$--$60.35$ and WikiText perplexity falls from $16.15$ to
$15.02$--$15.40$ --- but they no longer separate from each other: the three averages span
$0.29$ points ($60.06$ / $60.28$ / $60.35$) against per-task standard errors of $0.5$--$1.4$ points,
and the per-column wins split three ways. So the benefit of Kalman gating over the fixed-gain rule
is large and reproducible at both scales, while the extra structure of the diagonal gain does not
yet pay for itself over a single scalar on this suite. The $750$M and $1.3$B blocks disagree on
which KDN leads, which is the expected signature of a difference inside noise rather than a trend.

\begin{table}[t]
\centering
\small
\setlength{\tabcolsep}{4.5pt}
\begin{tabular}{@{}l cc cccccc c@{}}
\toprule
& \multicolumn{2}{c}{Perplexity $\downarrow$} & \multicolumn{6}{c}{Zero-shot accuracy (\%) $\uparrow$} & \\
\cmidrule(lr){2-3}\cmidrule(lr){4-9}
Model & Wiki. & LMB. & LMB. & PIQA & Hella. & Wino. & ARC-e & ARC-c & Avg. \\
\midrule
\multicolumn{10}{@{}l}{\emph{220M parameters, 20B tokens}}\\
DeltaNet        & 31.26 & 54.07 & 27.36 & 63.38 & 36.06 & 51.38 & 53.11 & 26.45 & 42.96 \\
Gated DeltaNet  & 30.83 & 46.89 & 27.34 & \underline{65.07} & 37.12 & 49.96 & 55.01 & 26.54 & 43.51 \\
GDN-2           & 31.62 & 42.95 & 30.18 & 63.98 & 35.47 & 51.62 & 54.17 & 26.28 & 43.62 \\
KDA             & 29.31 & \underline{36.39} & \underline{31.61} & 64.96 & \textbf{38.37} & \textbf{53.59} & 56.61 & \underline{27.39} & 45.42 \\
Diagonal KDN    & \underline{28.92} & 38.44 & 31.22 & \textbf{65.94} & \underline{37.96} & 51.30 & \textbf{58.54} & \textbf{27.73} & \underline{45.45} \\
Isotropic KDN   & \textbf{28.67} & \textbf{35.40} & \textbf{32.99} & \textbf{65.94} & 37.89 & \underline{52.49} & \underline{57.45} & 26.37 & \textbf{45.52} \\
\midrule
\multicolumn{10}{@{}l}{\emph{750M parameters, 50B tokens}}\\
GDN-2           & 21.20 & 17.88 & 41.18 & 70.08 & 46.90 & 54.54 & 64.27 & 31.74 & 51.45 \\
KDA             & 18.85 & 15.06 & \underline{44.34} & \textbf{70.95} & 51.23 & 54.93 & 67.72 & 34.04 & 53.87 \\
Isotropic KDN   & \textbf{18.42} & \underline{14.68} & 44.05 & \textbf{70.95} & \textbf{52.12} & \underline{56.12} & \underline{68.01} & \underline{35.24} & \underline{54.41} \\
Diagonal KDN    & \underline{18.64} & \textbf{14.15} & \textbf{45.66} & \textbf{70.95} & \underline{51.91} & \textbf{57.70} & \textbf{68.10} & \textbf{35.49} & \textbf{54.97} \\
\midrule
\multicolumn{10}{@{}l}{\emph{1.3B parameters, 100B tokens}}\\
GDN-2           & 16.15 & 11.29 & 49.45 & 72.63 & 57.80 & 59.19 & 72.73 & 39.25 & 58.51 \\
KDA             & 15.40 & \underline{10.09} & \textbf{51.15} & \textbf{74.16} & \textbf{60.61} & \underline{60.54} & \underline{73.48} & \underline{41.72} & \underline{60.28} \\
Isotropic KDN   & \underline{15.30} & \textbf{9.98} & \underline{50.96} & 73.23 & 60.23 & \textbf{61.33} & \textbf{74.71} & 41.64 & \textbf{60.35} \\
Diagonal KDN    & \textbf{15.02} & 10.60 & 49.91 & \underline{73.88} & \underline{60.29} & 60.14 & 73.32 & \textbf{42.83} & 60.06 \\
\bottomrule
\end{tabular}
\caption{Language modeling and zero-shot commonsense reasoning. Perplexity on WikiText and LAMBADA
(lower is better) and zero-shot accuracy (\%) on LAMBADA and the commonsense suite; \emph{Avg.} is
the mean of the six accuracies. All mixers are matched in parameter count and recurrent-state size
at each scale. Best per column within each block in \textbf{bold}, second best \underline{underlined}.}
\label{tab:pretrain-lm}
\end{table}

\paragraph{In-context retrieval.}
Table~\ref{tab:pretrain-niah} reports RULER single- and multi-key needle-in-a-haystack accuracy
across context lengths from 1K to 8K~\citep{hsieh2024ruler} (8K exceeds the 4K training length, i.e.
length extrapolation). We correct two issues in the released generator: essay contexts are fitted in
$25$-word rather than $500$-word increments, and each sample uses a deterministic random window from
the essay corpus instead of its fixed opening prefix. We also seed all keys, values, and needle depths,
so every model is evaluated on the same $500$ samples per cell.

At $220$M/$20$B retrieval is weak outside the easiest synthetic task. DeltaNet retains S-NIAH-1
farthest, while KDA leads S-NIAH-2 at 1K--2K; all models are near zero on long S-NIAH-3. At
$750$M/$50$B the separation is clear: the Diagonal KDN wins every nontrivial cell, remains perfect on
S-NIAH-1 through 8K, and reaches $71.2$ on S-NIAH-2 @ 4K and $71.0$ on S-NIAH-3 @ 2K. Its 14-cell
average is $70.3$, versus $62.4$ for the Isotropic KDN, $62.1$ for KDA, and $52.5$ for GDN-2.
This is consistent with an anisotropic gain protecting stored associations under interference.

At $1.3$B/$100$B the picture is mixed. GDN-2 and the Diagonal KDN tie on aggregate ($71.0$ and
$70.6$); the Diagonal KDN leads long S-NIAH-2, KDA leads S-NIAH-3, and GDN-2 leads two of three
multi-key lengths. KDA again degrades sharply on the easiest S-NIAH-1 probe. The previously observed
short-context collapse of the 1.3B Diagonal KDN disappears under randomized essay windows
(S-NIAH-2 @ 1K: $21.4\to88.0$; S-NIAH-3 @ 1K: $3.6\to82.4$), showing that it was an interaction
with the fixed opening essay rather than a short-memory failure. These are single checkpoints and one
window seed, so small differences should be read as directional.

\begin{table}[t]
\centering
\small
\setlength{\tabcolsep}{3pt}
\resizebox{\textwidth}{!}{%
\begin{tabular}{@{}l cccc cccc ccc ccc@{}}
\toprule
& \multicolumn{4}{c}{S-NIAH-1} & \multicolumn{4}{c}{S-NIAH-2} & \multicolumn{3}{c}{S-NIAH-3} & \multicolumn{3}{c}{MK-NIAH-1} \\
\cmidrule(lr){2-5}\cmidrule(lr){6-9}\cmidrule(lr){10-12}\cmidrule(lr){13-15}
Model & 1K & 2K & 4K & 8K & 1K & 2K & 4K & 8K & 1K & 2K & 4K & 1K & 2K & 4K \\
\midrule
\multicolumn{15}{@{}l}{\emph{220M parameters, 20B tokens}}\\
DeltaNet       & \underline{99.8} & \textbf{100.0} & \textbf{99.4} & \textbf{95.8} & 56.0 & 19.2 & 8.2 & 2.4 & \textbf{14.0} & \textbf{3.8} & \textbf{0.8} & \textbf{26.0} & \textbf{22.8} & \underline{12.0} \\
Gated DeltaNet & \textbf{100.0} & \underline{99.2} & \underline{92.8} & \underline{48.4} & 66.8 & 15.4 & 3.6 & 1.8 & 0.8 & 0.8 & 0.2 & 20.8 & 14.8 & 8.6 \\
GDN-2          & 91.4 & 48.0 & 13.2 & 3.8 & 59.4 & 19.6 & 7.0 & 3.0 & 1.4 & 0.2 & 0.0 & 20.6 & 18.4 & 11.4 \\
KDA            & \underline{99.8} & 97.2 & 68.8 & 22.8 & \textbf{84.8} & \textbf{35.0} & 7.0 & 2.0 & \underline{10.4} & \underline{1.2} & \textbf{0.8} & 22.2 & \underline{21.6} & \textbf{16.6} \\
Diagonal KDN   & 97.6 & 33.6 & 11.2 & 2.6 & \underline{69.8} & \underline{26.8} & \textbf{11.0} & \textbf{4.0} & 1.4 & 0.2 & \underline{0.4} & \underline{22.4} & 18.4 & 10.4 \\
Isotropic KDN  & 81.4 & 19.4 & 13.2 & 6.4 & 56.0 & 23.0 & \underline{9.2} & \underline{3.4} & 2.4 & 0.2 & \underline{0.4} & 19.8 & 17.8 & 11.2 \\
\midrule
\multicolumn{15}{@{}l}{\emph{750M parameters, 50B tokens}}\\
GDN-2         & \textbf{100.0} & \underline{99.8} & 90.2 & 47.2 & 95.0 & 74.6 & 36.0 & 12.6 & 64.6 & 31.0 & 9.2 & 27.8 & 25.6 & 22.0 \\
KDA           & \textbf{100.0} & \underline{99.8} & \underline{99.2} & 82.6 & \underline{96.4} & \underline{89.8} & \underline{63.2} & \underline{26.0} & 68.8 & 43.6 & 13.6 & \underline{32.6} & \underline{29.8} & \underline{24.2} \\
Isotropic KDN & \textbf{100.0} & \textbf{100.0} & 98.2 & \underline{89.4} & 96.0 & 86.0 & 47.8 & 16.2 & \underline{80.4} & \underline{58.4} & \underline{22.6} & 30.4 & 25.4 & 22.2 \\
Diagonal KDN  & \textbf{100.0} & \textbf{100.0} & \textbf{100.0} & \textbf{100.0} & \textbf{98.6} & \textbf{91.4} & \textbf{71.2} & \textbf{29.0} & \textbf{84.2} & \textbf{71.0} & \textbf{32.2} & \textbf{41.0} & \textbf{36.4} & \textbf{29.6} \\
\midrule
\multicolumn{15}{@{}l}{\emph{1.3B parameters, 100B tokens}}\\
GDN-2         & \textbf{100.0} & \textbf{100.0} & \textbf{100.0} & \textbf{99.8} & \textbf{98.6} & \underline{92.4} & 67.4 & \underline{28.0} & \underline{88.0} & 64.4 & 29.4 & \underline{50.6} & \textbf{44.0} & \textbf{31.4} \\
KDA           & \underline{99.8} & \underline{99.2} & \underline{66.6} & 30.2 & 98.0 & \textbf{95.0} & 62.2 & 23.6 & \textbf{92.6} & \textbf{77.2} & \textbf{43.2} & 47.4 & 38.8 & 26.6 \\
Isotropic KDN & \textbf{100.0} & \textbf{100.0} & \textbf{100.0} & 94.6 & \underline{98.4} & 91.8 & \underline{68.4} & 24.2 & 86.8 & \underline{67.6} & 34.6 & 41.6 & 29.8 & 23.4 \\
Diagonal KDN  & \textbf{100.0} & \textbf{100.0} & \textbf{100.0} & \underline{99.6} & 88.0 & 91.0 & \textbf{77.2} & \textbf{29.4} & 82.4 & 65.0 & \underline{36.0} & \textbf{51.0} & \underline{42.2} & \underline{27.2} \\
\bottomrule
\end{tabular}%
}
\caption{In-context retrieval accuracy (\%) on RULER single- and multi-key needle-in-a-haystack
tasks~\citep{hsieh2024ruler} across context lengths 1K--8K. All models are recurrent-only and matched
in parameter count and recurrent-state size at each scale; 8K exceeds the 4K training length (length
extrapolation). Essay contexts use deterministic random corpus windows and a 25-word fitting increment;
all models receive the same 500 samples per cell. Best per column within each block in \textbf{bold},
second best \underline{underlined}.}
\label{tab:pretrain-niah}
\end{table}


% =============================================================================
\section{Related Work}
\label{sec:related-work}

\paragraph{Linear attention, delta rules, and gated recurrent mixers.}
Linear attention replaces the quadratic softmax attention matrix with a recurrent fast-weight
state, giving constant-memory autoregressive inference and scan-parallel training
\citep{katharopoulos2020transformers}. DeltaNet interprets this state as an online associative
memory and updates it by writing only the value residual at the current key
\citep{schlag2021linear}. Later gated variants add learned forgetting and more expressive
state transitions, including Gated Linear Attention, Gated DeltaNet, and KDA
\citep{yang2023gated,yang2024gdn,kimiteam2025kimilinear}. Our starting point is this
associative-memory view: the recurrent state is a key--value map $\ve{S}_t$, the token
measurement is $\ve{v}_t\approx \ve{S}_t^{\tr}\ve{k}_t$, and the write is a Kalman residual
update. This lets us place DeltaNet, GDN, and KDA inside a single state-space filtering
picture as fixed-gain or restricted-covariance approximations.

\paragraph{Kalman filtering as a sequence mixer.}
Classical Kalman filtering performs optimal online inference in linear--Gaussian state-space
models by maintaining both a posterior mean and covariance. Recent work by
\citet{shaj2026kalmanlinearattention} also introduces a method named Kalman Linear Attention
and shows that a diagonal information-form Kalman filter can be made scan-parallel through a
M\"obius precision recursion. Their layer treats token values as noisy observations of a
per-channel latent feature state, uses an OU-discretized stochastic prior with time-invariant
dynamics, and reads out the posterior mean by an elementwise query projection. In contrast, our
KDN studies Kalman filtering for \emph{delta-network associative memory}: the latent object is
the memory matrix $\ve{S}_t\in\R^{d_k\times d_v}$, the observation is a key-conditioned value
measurement $\ve{v}_t=\ve{S}_t^{\tr}\ve{k}_t+\ve{e}_t$, and the central object is the Kalman
write direction $\ve{\kappa}_t$. The exact filter for this model carries a dense key-space
covariance, and our diagonal and isotropic variants are structured approximations to that exact
associative-memory update. Thus the two papers share information-form scan algebra, but they
apply it to different latent states and different notions of attention.

\paragraph{Parallel scans and information-form recurrences.}
The scanability of KDN comes from representing uncertainty recursions as associative maps.
For our diagonal construction, the additive process-noise prediction followed by the diagonal
information update yields a per-channel M\"obius map, so prefix uncertainty states are computed by
products of $2\times2$ matrices. Once the gains are available, the memory update is a standard
delta-family affine recurrence and can reuse the same scan machinery developed for linear
attention and DeltaNet \citep{yang2024parallelizing}. This separates our contribution into two
parts: a Kalman-derived gain scan, and a conventional fast-weight memory scan conditioned on
those gains.

\paragraph{Positioning.}
Our goal is not only to add uncertainty to a recurrent mixer, but to identify which parts of
Kalman filtering are compatible with key--value recall. Exact dense Kalman is the filtering
reference for the associative-memory model, but its covariance-optimal gain can rotate writes
away from the query key. The isotropic and diagonal variants therefore expose a useful design
axis: scalar uncertainty gives an adaptive write strength while preserving the delta-rule
binding direction, whereas anisotropic uncertainty gives a richer but potentially less
retrieval-aligned write direction. This distinction is specific to the associative-memory
formulation and is complementary to prior work that uses Kalman filtering primarily as a
per-channel SSM primitive.


% =============================================================================
\section{Future Directions}
\label{sec:future-directions}

The diagonal-information construction is the conservative path. The same state-space view
also suggests several extensions that trade additional structure for a closer approximation
to the Kalman update.

\paragraph{Direction A: richer process models.}
\begin{enumerate}[leftmargin=1.8em,itemsep=3pt,topsep=2pt,label=\textbf{A\arabic*.}]
  \item \textbf{Block-diagonal transitions.} Use
  $\ve{D}_t=\operatorname{blockdiag}(\ve{D}_t^{(1)},\ldots,\ve{D}_t^{(m)})$. Small blocks allow local
  mixing among related memory channels while keeping the transition structured enough for
  parallel composition.

  \item \textbf{Basis-rotated diagonal transitions.} Use
  $\ve{D}_t=\ve{R}\diag(\ve{\alpha}_t)\ve{R}^{\tr}$ with a fixed or slowly changing basis $\ve{R}$. Forgetting then
  acts in learned memory directions rather than raw coordinates, but the per-token decay
  remains diagonal in that basis.

  \item \textbf{Prefix-conditioned transitions.} Let $\ve{D}_t$ depend on scan-computable
  statistics $\ve{z}_t$, such as decayed key energy or an innovation proxy, instead of depending
  directly on the sequential state $\ve{S}_{t-1}$. Then $\ve{D}_t=\ve{D}(\ve{x}_t,\ve{z}_t)$ can react to drift while
  keeping $\ve{A}_t$ available before the scan.
\end{enumerate}

\paragraph{Direction B: richer covariance and gain families.}
\begin{enumerate}[leftmargin=1.8em,itemsep=3pt,topsep=2pt,label=\textbf{B\arabic*.}]
  \item \textbf{Transition-aware scalar covariance.} Set $\ve{P}_t=c_t\ve{I}$, but predict the scalar
  uncertainty through $\ve{D}_t$ before forming the gain. This still recovers the delta-rule
  direction $\ve{\kappa}_t=\beta_t \ve{k}_t$, while allowing the scalar write strength to depend on
  how strongly the memory has just been forgotten.

  \item \textbf{Block-diagonal covariance.} Track $\ve{P}_t$ or $\ve{C}_t=\ve{P}_t^{-1}$ in small blocks.
  This captures correlations among nearby key channels while replacing one large inverse by
  many small inverses. The prediction remains local:
  $\widehat{\ve{P}}_t^{(j)}=\ve{D}_t^{(j)}\ve{P}_{t-1}^{(j)}\ve{D}_t^{(j)\tr}+\ve{\Omega}_t^{(j)}$.

  \item \textbf{Diagonal-plus-low-rank covariance.} Use
  $\ve{P}_t=\diag(\ve{p}_t)+\ve{U}_t\ve{U}_t^{\tr}$, with $\ve{p}_t$ and $\ve{U}_t$ produced from prefix statistics.
  Woodbury-style reads can keep the gain tractable, while the low-rank term captures global
  uncertainty directions missed by a diagonal approximation.

  \item \textbf{Chunk-boundary covariance updates.} Update $\ve{P}_t$ only at chunk boundaries,
  while freezing the gain inside each chunk. This gives a TTT-like approximation: the
  expensive uncertainty update is carried by a short sequential pass across chunks, but all
  token updates within a chunk remain fixed-gain affine scans.
\end{enumerate}


\bibliography{main}
\bibliographystyle{plainnat}


\appendix

% =============================================================================
\section{Theoretical Proofs}
\label{app:proofs}

\begin{proof}[Proof of Proposition~\ref{prop:kalman-optimal-update} (Kalman Optimal Update)]

Let the associative-memory state be
\[
  \ve{S}_t \in \R^{d_k\times d_v},
  \qquad
  \mathcal{M}_t(\ve{k})=\ve{S}_t^{\tr}\ve{k}.
\]
The linear--Gaussian state-space model is
\begin{align}
  \ve{S}_t^\star
  &= \ve{D}_t\ve{S}_{t-1}^\star+\ve{W}_t,
  &
  \ve{W}_t&\sim\mathcal{N}(0,\ve{\Omega}_t), \\
  \ve{v}_t
  &= \ve{S}_t^{\star\tr}\ve{k}_t+\ve{e}_t,
  &
  \ve{e}_t&\sim\mathcal{N}(0,r_t\ve{I}).
\end{align}
Let $\ve{z}_s=(\ve{k}_s,\ve{v}_s)$ denote the key--value observation at step $s$.
At the end of step $t-1$, the Kalman filter represents the posterior over the latent
memory as a Gaussian,
\begin{equation}
  p(\ve{S}_{t-1}^\star\mid \ve{z}_{1:t-1})
  =
  \mathcal{N}(\ve{S}_{t-1},\ve{P}_{t-1}),
  \label{eq:previous-filtering-posterior}
\end{equation}
where $\ve{S}_{t-1}$ is the posterior mean, i.e.\ the point estimate of the memory after
processing the first $t-1$ key--value pairs, and $\ve{P}_{t-1}$ is its key-space uncertainty.
For the matrix memory, this covariance is understood per value coordinate, with the same
key-space covariance shared across value columns.

Before observing token $t$, the filter first predicts the next latent memory by marginalizing
over the previous posterior:
\begin{equation}
  p(\ve{S}_t^\star\mid \ve{z}_{1:t-1})
  =
  \int
  p(\ve{S}_t^\star\mid \ve{S}_{t-1}^\star)
  p(\ve{S}_{t-1}^\star\mid \ve{z}_{1:t-1})
  \,d\ve{S}_{t-1}^\star .
  \label{eq:prediction-integral}
\end{equation}
Because both factors are linear--Gaussian, the predicted distribution is Gaussian,
\begin{equation}
  p(\ve{S}_t^\star\mid \ve{z}_{1:t-1})
  =
  \mathcal{N}(\widehat{\ve{S}}_t,\widehat{\ve{P}}_t).
  \label{eq:predicted-prior}
\end{equation}
Its mean and covariance are the exact Kalman prediction:
\begin{equation}
  \widehat{\ve{S}}_t=\ve{D}_t\ve{S}_{t-1},
  \qquad
  \widehat{\ve{P}}_t
  =
  \ve{D}_t\ve{P}_{t-1}\ve{D}_t^{\tr}+\ve{\Omega}_t.
  \label{eq:kam-predict}
\end{equation}
Thus $\widehat{\ve{S}}_t$ is a point estimate, but specifically the \emph{predicted prior
mean} of $\ve{S}_t^\star$ conditioned on the past key--value sequence
$\ve{z}_{1:t-1}$. It is not yet the point estimate of
$p(\ve{S}_t^\star\mid \ve{z}_{1:t})$, because the current pair
$(\ve{k}_t,\ve{v}_t)$ has not been used. The covariance $\widehat{\ve{P}}_t$ is not a point
estimate; it is the predicted uncertainty around that mean.
The predicted read and innovation are
\begin{equation}
  \widehat{\ve{v}}_t=\widehat{\ve{S}}_t^{\tr}\ve{k}_t,
  \qquad
  \ve{\delta}_t=\ve{v}_t-\widehat{\ve{v}}_t.
  \label{eq:kam-innovation}
\end{equation}

This residual form is the matrix version of the standard Kalman correction. To see it,
first look at one value coordinate $j$. Let $\ve{s}_{t,j}$ be column $j$ of the memory and
let $v_{t,j}$ be coordinate $j$ of the observed value. The measurement model is
\begin{equation}
  v_{t,j}
  =
  \ve{k}_t^{\tr}\ve{s}_{t,j}^\star+e_{t,j},
  \qquad
  e_{t,j}\sim\mathcal{N}(0,r_t).
  \label{eq:single-value-coordinate-observation}
\end{equation}
Before this measurement is used, the predicted distribution for this column is
\begin{equation}
  \ve{s}_{t,j}^\star \mid \ve{z}_{1:t-1}
  \sim
  \mathcal{N}(\widehat{\ve{s}}_{t,j},\widehat{\ve{P}}_t).
  \label{eq:single-coordinate-prior}
\end{equation}
The observation $v_{t,j}$ is one scalar. Its predicted mean is
$\ve{k}_t^{\tr}\widehat{\ve{s}}_{t,j}$, so the part of the observation not already predicted
by the prior is the scalar residual
\begin{equation}
  \delta_{t,j}
  =
  v_{t,j}-\ve{k}_t^{\tr}\widehat{\ve{s}}_{t,j}.
  \label{eq:single-coordinate-innovation}
\end{equation}
For a joint Gaussian pair $(x,y)$ with scalar $y$, the conditional mean is
\[
  \mathbb{E}[x\mid y]
  =
  \mathbb{E}[x]
  +
  \operatorname{Cov}(x,y)\operatorname{Var}(y)^{-1}
  \big(y-\mathbb{E}[y]\big).
\]
Apply this identity with $x=\ve{s}_{t,j}^\star$ and $y=v_{t,j}$. Under
\eqref{eq:single-coordinate-prior}--\eqref{eq:single-value-coordinate-observation},
\begin{equation}
  \operatorname{Cov}(\ve{s}_{t,j}^\star,v_{t,j})
  =
  \widehat{\ve{P}}_t\ve{k}_t,
  \qquad
  \operatorname{Var}(v_{t,j})
  =
  \ve{k}_t^{\tr}\widehat{\ve{P}}_t\ve{k}_t+r_t.
  \label{eq:single-coordinate-cross-cov}
\end{equation}
Therefore the posterior mean is
\begin{equation}
  \ve{s}_{t,j}
  =
  \widehat{\ve{s}}_{t,j}
  +
  \ve{\kappa}_t\,\delta_{t,j}.
  \label{eq:single-coordinate-kalman-correction}
\end{equation}
where
\begin{equation}
  \ve{\kappa}_t
  =
  \frac{\widehat{\ve{P}}_t\ve{k}_t}
  {\ve{k}_t^{\tr}\widehat{\ve{P}}_t\ve{k}_t+r_t}.
  \label{eq:single-coordinate-gain}
\end{equation}
The same $\ve{\kappa}_t$ applies to every value coordinate because the model uses the same
key-space covariance $\widehat{\ve{P}}_t$ and the same observation noise $r_t$ for every
coordinate of $\ve{v}_t$. Stacking \eqref{eq:single-coordinate-kalman-correction} over
$j=1,\ldots,d_v$ gives the rank-one residual write
\begin{equation}
  \ve{S}_t
  =
  \widehat{\ve{S}}_t
  +
  \ve{\kappa}_t
  \big(\ve{v}_t-\widehat{\ve{S}}_t^{\tr}\ve{k}_t\big)^{\tr}.
  \label{eq:kam-residual-write}
\end{equation}
Thus the residual write does not come from DeltaNet first; it comes from the Kalman
posterior-mean update $x^+=x^-+K(y-Hx^-)$. DeltaNet later appears as the fixed-gain special
case where $\ve{\kappa}_t$ is forced to be proportional to $\ve{k}_t$.
For a candidate write direction $\ve{\kappa}_t$, the posterior covariance induced by this
write is
\begin{equation}
  \ve{P}_t(\ve{\kappa}_t)
  =
  (\ve{I}-\ve{\kappa}_t\ve{k}_t^{\tr})
  \widehat{\ve{P}}_t
  (\ve{I}-\ve{\kappa}_t\ve{k}_t^{\tr})^{\tr}
  +
  r_t\ve{\kappa}_t\ve{\kappa}_t^{\tr}.
  \label{eq:kam-candidate-cov}
\end{equation}
The Kalman write direction is the minimizer of total posterior uncertainty,
\begin{equation}
  \ve{\kappa}_t^\star
  =
  \arg\min_{\ve{\kappa}_t}
  \Trace\!\left(\ve{P}_t(\ve{\kappa}_t)\right).
  \label{eq:kam-trace-objective}
\end{equation}
Expanding the trace gives
\begin{equation}
  \Trace\!\left(\ve{P}_t(\ve{\kappa}_t)\right)
  =
  \Trace(\widehat{\ve{P}}_t)
  -2\ve{\kappa}_t^{\tr}\widehat{\ve{P}}_t\ve{k}_t
  +
  \big(r_t+\ve{k}_t^{\tr}\widehat{\ve{P}}_t\ve{k}_t\big)
  \ve{\kappa}_t^{\tr}\ve{\kappa}_t .
  \label{eq:kam-trace-expanded}
\end{equation}
Taking the gradient with respect to $\ve{\kappa}_t$ and setting it to zero yields
\begin{equation}
  \boxed{
  \ve{\kappa}_t
  =
  \frac{\widehat{\ve{P}}_t\ve{k}_t}
  {r_t+\ve{k}_t^{\tr}\widehat{\ve{P}}_t\ve{k}_t}
  }.
  \label{eq:kam-optimal-gain}
\end{equation}
Therefore the exact Kalman Associative Memory update is
\begin{equation}
\begin{aligned}
  \widehat{\ve{S}}_t
  &=\ve{D}_t\ve{S}_{t-1},
  &
  \widehat{\ve{P}}_t
  &=\ve{D}_t\ve{P}_{t-1}\ve{D}_t^{\tr}+\ve{\Omega}_t,
  \\
  \ve{\kappa}_t
  &=
  \frac{\widehat{\ve{P}}_t\ve{k}_t}
  {r_t+\ve{k}_t^{\tr}\widehat{\ve{P}}_t\ve{k}_t},
  &
  \ve{S}_t
  &=
  \widehat{\ve{S}}_t
  +
  \ve{\kappa}_t
  \big(\ve{v}_t-\widehat{\ve{S}}_t^{\tr}\ve{k}_t\big)^{\tr},
  \\
  \ve{P}_t
  &=
  (\ve{I}-\ve{\kappa}_t\ve{k}_t^{\tr})\widehat{\ve{P}}_t.
\end{aligned}
\label{eq:kam-optimal-update}
\end{equation}
This is the reference update. The remaining question is which part of it can be made
tractable and scan-computable without losing the link between uncertainty, transition, and
write strength.
\end{proof}

% =============================================================================
\section{Background: the Kalman filter}
\label{app:kalman}

The Kalman filter is the optimal recursive estimator for a linear--Gaussian
state-space model; \S\ref{sec:memory-kalman} casts the linear-attention memory as one
instance. We recall the classical form and its specializations here.

\paragraph{Linear--Gaussian state-space model.}
A hidden state $\ve{x}_t\in\R^n$ evolves and is observed through
\begin{equation}
\begin{aligned}
  \ve{x}_t &= \ve{F}_t \ve{x}_{t-1} + \ve{w}_t,
  & \ve{w}_t &\sim \N(0,\ve{\Omega}_t)\quad\text{(process / dynamics)},\\[2pt]
  \ve{y}_t &= \ve{H}_t \ve{x}_t + \ve{e}_t,
  & \ve{e}_t &\sim \N(0,\ve{R}_t)\quad\text{(observation)},
\end{aligned}
\label{eq:kf-ssm}
\end{equation}
with transition $\ve{F}_t$, observation map $\ve{H}_t$, and independent zero-mean Gaussian
process and measurement noise of covariance $\ve{\Omega}_t$ and $\ve{R}_t$. Given $\ve{y}_{1:t}$ the
posterior over the state is Gaussian, $p(\ve{x}_t\mid \ve{y}_{1:t})=\N(\widehat{\ve{x}}_t,\ve{P}_t)$, and the
filter propagates its mean $\widehat{\ve{x}}_t$ and covariance $\ve{P}_t$.

\paragraph{The filtering recursion.}
Each step is a \emph{predict} (time update) then an \emph{update} (measurement update):
\begin{equation}
\begin{aligned}
  \text{predict:}\quad & \widehat{\ve{x}}_t^- = \ve{F}_t \widehat{\ve{x}}_{t-1},
  & \ve{P}_t^- &= \ve{F}_t \ve{P}_{t-1} \ve{F}_t^{\tr} + \ve{\Omega}_t,\\[2pt]
  \text{update:}\quad  & \widehat{\ve{x}}_t = \widehat{\ve{x}}_t^- + \ve{K}_t\,\ve{\nu}_t,
  & \ve{P}_t &= (\ve{I} - \ve{K}_t \ve{H}_t)\,\ve{P}_t^-,
\end{aligned}
\label{eq:kf-rec}
\end{equation}
driven by the \emph{innovation} (prediction error) and \emph{Kalman gain}
\begin{equation}
  \ve{\nu}_t = \ve{y}_t - \ve{H}_t \widehat{\ve{x}}_t^-, \qquad
  \ve{K}_t = \ve{P}_t^- \ve{H}_t^{\tr}\,\ve{\Sigma}_t^{-1}, \qquad
  \ve{\Sigma}_t = \ve{H}_t \ve{P}_t^- \ve{H}_t^{\tr} + \ve{R}_t ,
  \label{eq:kf-gain}
\end{equation}
where $\ve{\Sigma}_t$ is the innovation covariance. The gain weighs new evidence against
the prior by relative uncertainty: $\ve{K}_t$ is large when the prior is uncertain ($\ve{P}_t^-$
large) or the measurement reliable ($\ve{R}_t$ small), and small otherwise.

\paragraph{Optimality.}
For a linear--Gaussian model \eqref{eq:kf-ssm} the filter is exact: $\widehat{\ve{x}}_t$ is the
posterior mean, simultaneously the minimum-mean-square-error and maximum-a-posteriori
estimate. Without Gaussianity it remains the best \emph{linear} unbiased estimator. It
is thus the optimal way to fuse a drifting prior with streaming observations.

\paragraph{Special case: recursive least squares.}
Take a \emph{static} state ($\ve{F}_t=\ve{I}$, $\ve{\Omega}_t=0$) observed through a regressor,
$y_t=\ve{\phi}_t^{\tr}\ve{\theta}+e_t$ (so $\ve{H}_t=\ve{\phi}_t^{\tr}$, scalar $R_t=r$). Then
\eqref{eq:kf-rec}--\eqref{eq:kf-gain} reduce to \emph{recursive least squares}, which
minimizes $\sum_{\tau\le t}\|y_\tau-\ve{\phi}_\tau^{\tr}\ve{\theta}\|_2^2$ online; the covariance
$\ve{P}_t=\big(\sum_{\tau\le t}\ve{\phi}_\tau\ve{\phi}_\tau^{\tr}\big)^{-1}$ is the inverse data Gram
matrix. An exponential \emph{forgetting factor} $\lambda\in(0,1]$ --- weighting old
data by $\lambda^{t-\tau}$ --- is exactly the Kalman filter with process noise
$\ve{\Omega}_t=(\lambda^{-1}-1)\ve{P}_{t-1}$ (equivalently covariance inflation $\ve{P}_t^-=\ve{P}_{t-1}/\lambda$),
letting the estimate track slowly drifting parameters. RLS is therefore the
static-state (or random-walk) special case of the Kalman filter.

% =============================================================================
\subsection{The associative-memory Kalman filter}
\label{app:mem-kf}

\paragraph{State-space model.}
Write the memory itself as the hidden state: it drifts between tokens and is observed, one
noisy key--value pair at a time, through the current key,
\begin{equation}
\begin{aligned}
  \ve{S}_t &= \ve{D}_t \ve{S}_{t-1} + \ve{W}_t, & \ve{W}_t &\sim\N(0,\ve{\Omega}_t),\\[2pt]
  \ve{v}_t &= \ve{S}_t^{\tr}\ve{k}_t + \ve{e}_t, & \ve{e}_t &\sim\N(0,r_t),
\end{aligned}
\label{eq:mem-ssm}
\end{equation}
with transition/decay $\ve{D}_t\in\R^{d_k\times d_k}$, drift covariance $\ve{\Omega}_t$, and scalar
observation noise $r_t$. The observation map $\ve{k}_t^{\tr}$ and the dynamics $\ve{D}_t$ are shared by
all $d_v$ value channels, so the posterior factorizes across them: we track one estimate
$\ve{S}_t\in\R^{d_k\times d_v}$ and a single key-space covariance $\ve{P}_t\in\R^{d_k\times d_k}$.

\paragraph{Optimal solution.}
The minimum-variance estimate is the Kalman recursion \eqref{eq:kf-rec}--\eqref{eq:kf-gain}
specialized to \eqref{eq:mem-ssm}:
\begin{equation}
\begin{aligned}
  \text{predict:}\quad & \widehat{\ve{S}}_t = \ve{D}_t \ve{S}_{t-1},
  & \widehat{\ve{P}}_t &= \ve{D}_t \ve{P}_{t-1} \ve{D}_t^{\tr} + \ve{\Omega}_t,\\[2pt]
  \text{update:}\quad  & \ve{S}_t = \widehat{\ve{S}}_t
  + \ve{\kappa}_t\,(\ve{v}_t - \widehat{\ve{S}}_t^{\tr}\ve{k}_t)^{\tr},
  & \ve{P}_t &= (\ve{I} - \ve{\kappa}_t \ve{k}_t^{\tr})\widehat{\ve{P}}_t,
\end{aligned}
\label{eq:mem-kf}
\end{equation}
with gain $\ve{\kappa}_t = \widehat{\ve{P}}_t \ve{k}_t/(r_t + \ve{k}_t^{\tr}\widehat{\ve{P}}_t \ve{k}_t)$. In the static-state limit
($\ve{D}_t=\ve{I}$, $\ve{\Omega}_t=0$) the estimate has the closed weighted-least-squares form
\begin{equation}
  \ve{S}_t = \ve{P}_t \sum_{\tau\le t}\gamma_{t,\tau}\,\frac{1}{r_\tau}\ve{k}_\tau \ve{v}_\tau^{\tr},
  \qquad
  \ve{P}_t = \Big(\sum_{\tau\le t}\gamma_{t,\tau}\,\frac{1}{r_\tau}
  \ve{k}_\tau \ve{k}_\tau^{\tr}+\ve{P}_0^{-1}\Big)^{-1},
  \label{eq:mem-wls}
\end{equation}
which \eqref{eq:mem-kf} computes online; the gain $\ve{\kappa}_t$ is a data- and state-dependent
write strength that replaces the delta rule's fixed $\beta_t \ve{k}_t$.

\paragraph{Why it resists parallelization.}
The delta-rule family trains in parallel because its per-token update is an \emph{affine}
map of the state, $\ve{S}_t=\ve{A}_t \ve{S}_{t-1}+\ve{b}_t$, whose coefficients
$\ve{A}_t=(\ve{I}-\beta_t \ve{k}_t \ve{k}_t^{\tr})\ve{D}_t$ and
$\ve{b}_t=\beta_t \ve{k}_t \ve{v}_t^{\tr}$ depend on the \emph{inputs
alone}; a chunk of such maps composes into dense matmuls (the WY / chunked scan). The
optimal filter \eqref{eq:mem-kf} breaks this on two counts. First, the covariance obeys a
\emph{Riccati recursion}
$\ve{P}_t=(\ve{I}-\ve{\kappa}_t \ve{k}_t^{\tr})(\ve{D}_t \ve{P}_{t-1}\ve{D}_t^{\tr}+\ve{\Omega}_t)$ that is
nonlinear (rational) in $\ve{P}_{t-1}$ and has no closed-form chunk composition --- it must be
iterated token by token. Second, the gain $\ve{\kappa}_t$ depends on $\ve{P}_{t-1}$, so the state
coefficient $\ve{A}_t$ depends on the \emph{accumulated} covariance rather than the current
input, and the recurrence is no longer a fixed linear scan. (Letting $\ve{D}_t$ or $\ve{\Omega}_t$ depend
on $\ve{S}_{t-1}$ compounds this, making even the predict step nonlinear.) On top of the
sequential dependence, the covariance is $d_k\times d_k$: $O(d_k^2)$ state and work per
step, versus the delta family's $O(d_k d_v)$ matmuls.

\paragraph{Parallelizable relaxations.}
Parallel training needs the per-token coefficients to be either input-only or computable by
an \emph{associative scan} (a prefix operation). Several general relaxations recover this
while keeping some of the second-order/state dependence the delta family discards:
\begin{itemize}[leftmargin=1.4em,itemsep=3pt,topsep=2pt]
  \item \textbf{Structured covariance.} Restrict $\ve{P}_t$ to a cheap form. A diagonal
  $\ve{P}_t=\diag(\ve{p}_t)$ decouples the Riccati into $d_k$ scalar recursions ($O(d_k)$ rather than
  $O(d_k^2)$); a scalar $\ve{P}_t=b_t \ve{I}$ collapses it to one recursion and a per-step scalar gain
  $\beta_t=b_t/(r_t+b_t)$ for normalized keys; a diagonal-plus-low-rank $\ve{P}_t$ interpolates.
  \item \textbf{Gain from prefix statistics.} Replace the exact Riccati by a gain that is a
  pointwise function of \emph{cumulative} quantities a scan can produce --- e.g.\ a running
  key energy $\sum_{\tau\le t}\gamma_{t,\tau}\ve{k}_\tau \ve{k}_\tau^{\tr}$ (or its diagonal), or an
  EMA of the innovation. The prefix-sum is parallel; the nonlinearity is applied pointwise.
  \item \textbf{Chunk-wise freezing.} Hold the covariance/gain (and any state-dependent
  transition) fixed within a chunk at its chunk-start value, updating it only across chunks.
  Intra-chunk the update is the fixed-gain delta rule (parallel matmuls); inter-chunk a short
  sequential pass carries $\ve{P}$. This interpolates between the delta family (never update) and
  the exact filter (chunk size one), mirroring mini-batch test-time training.
  \item \textbf{Learned surrogate.} Replace the Riccati solution by a small learned map from
  cheap, parallelizable features (the input, prefix statistics, the innovation) to the gain
  and transition --- a parametric stand-in for the second-order recursion, trained end to
  end. State-dependent gates are one member of this class.
\end{itemize}
Each trades fidelity to the Kalman optimum for parallelism and cost; the delta family is the
extreme that freezes the covariance entirely, and richer choices move back toward the optimum
along one of these axes.


% =============================================================================
\subsection{Associative memory as recursive least squares}
\label{app:mem-rls}

\subsubsection{Static memory: recursive least squares}
\label{app:mem-rls-static}

\paragraph{Online least-squares objective.}
Dropping the drift ($\ve{D}_t=\ve{I}$, $\ve{\Omega}_t=0$) reduces the filter of \S\ref{app:mem-kf} to estimating
a \emph{fixed} memory from streaming observations --- ridge-regularized, recency-weighted
least squares,
\begin{equation}
  \ve{S}_t = \argmin_{\ve{S}}\ \tfrac12\sum_{\tau\le t}\lambda^{t-\tau}
  \big\|\ve{S}^{\tr}\ve{k}_\tau-\ve{v}_\tau\big\|_2^2
        + \tfrac12\lambda^{t}\mu\|\ve{S}\|_F^2 ,
  \label{eq:rls-obj}
\end{equation}
with forgetting factor $\lambda\in(0,1]$ and ridge $\mu>0$ (a decaying prior); this is the
static-state case of \eqref{eq:mem-wls}.

\paragraph{The RLS recursion.}
Zeroing the gradient gives the normal equations $\ve{C}_t \ve{S}_t = \ve{B}_t$ with the \emph{information
matrix} and target
\begin{equation}
  \ve{C}_t = \sum_{\tau\le t}\lambda^{t-\tau}\ve{k}_\tau \ve{k}_\tau^{\tr}
  + \lambda^{t}\mu \ve{I},
  \qquad
  \ve{B}_t = \sum_{\tau\le t}\lambda^{t-\tau}\ve{k}_\tau \ve{v}_\tau^{\tr},
  \label{eq:rls-info}
\end{equation}
both accumulating linearly, $\ve{C}_t=\lambda \ve{C}_{t-1}+\ve{k}_t \ve{k}_t^{\tr}$ and
$\ve{B}_t=\lambda \ve{B}_{t-1}+\ve{k}_t \ve{v}_t^{\tr}$. With $\ve{P}_t=\ve{C}_t^{-1}$, the Sherman--Morrison identity gives
the classical \emph{recursive least squares} update
\begin{equation}
\begin{aligned}
  \ve{g}_t &= \frac{\ve{P}_{t-1}\ve{k}_t}
  {\lambda+\ve{k}_t^{\tr}\ve{P}_{t-1}\ve{k}_t}, \\[3pt]
  \ve{S}_t &= \ve{S}_{t-1}+\ve{g}_t\,
  (\ve{v}_t-\ve{S}_{t-1}^{\tr}\ve{k}_t)^{\tr}, \\[3pt]
  \ve{P}_t &= \lambda^{-1}\big(\ve{P}_{t-1}-\ve{g}_t \ve{k}_t^{\tr}\ve{P}_{t-1}\big),
\end{aligned}
\label{eq:rls-rec}
\end{equation}
the exact minimizer of \eqref{eq:rls-obj} computed online.

\paragraph{RLS is online Newton.}
Since $\ve{g}_t=\ve{P}_t \ve{k}_t$, the update \eqref{eq:rls-rec} is a \emph{preconditioned} gradient step on
the per-token recall loss $\mc L_t(\ve{S})=\tfrac12\|\ve{S}^{\tr}\ve{k}_t-\ve{v}_t\|_2^2$,
\begin{equation}
  \ve{S}_t = \ve{S}_{t-1} - \ve{P}_t\,\nabla_{\!\ve{S}}\mc L_t(\ve{S}_{t-1}),
  \qquad \ve{P}_t = \ve{C}_t^{-1},
  \label{eq:rls-newton}
\end{equation}
whose preconditioner $\ve{P}_t$ is the inverse accumulated Hessian --- so RLS is \emph{online
Newton} (a natural-gradient step). The delta-rule family replaces $\ve{P}_t$ by a fixed
$\beta_t \ve{I}$: \emph{online SGD}, the first-order ($\ve{P}=\ve{I}$) approximation of this second-order
($\ve{P}=\ve{C}^{-1}$) optimum, with GDN/KDA adding an $\ell_2$ shrinkage (the decay) on top.

\paragraph{Information form: a linear accumulator.}
The obstruction to parallel training is \emph{not} the accumulation: by \eqref{eq:rls-info}
the information matrix $\ve{C}_t$ (and $\ve{B}_t$) is a decayed cumulative sum --- a prefix scan, hence
parallelizable. The nonlinearity is confined to the per-step inverse $\ve{P}_t=\ve{C}_t^{-1}$ (or
$\ve{S}_t=\ve{C}_t^{-1}\ve{B}_t$). Carrying $\ve{C}_t$ rather than $\ve{P}_t$ --- the \emph{information} form --- is
therefore the parallel-friendly representation: the state accumulates linearly, and only the
read-out inverts.

\subsubsection{Drifting memory: dynamic least squares}
\label{app:mem-rls-drift}

Adding the drift $\ve{D}_t$ changes the optimization problem qualitatively: the unknown is no
longer one fixed memory $\ve{S}$, but a trajectory of latent memories
$\ve{S}_{0:t}=(\ve{S}_0,\ldots,\ve{S}_t)$ coupled by the process model. The batch MAP problem
corresponding to \eqref{eq:mem-ssm} is
\begin{equation}
  \{\ve{S}_\tau^\star\}_{\tau=0}^t
  =
  \argmin_{\ve{S}_0,\ldots,\ve{S}_t}
  \frac12\|\ve{S}_0\|_{\ve{P}_0^{-1}}^2
  +\frac12\sum_{\tau=1}^t
  \|\ve{S}_\tau-\ve{D}_\tau\ve{S}_{\tau-1}\|_{\ve{\Omega}_\tau^{-1}}^2
  +\frac12\sum_{\tau=1}^t
  \frac{1}{r_\tau}
  \|\ve{S}_\tau^{\tr}\ve{k}_\tau-\ve{v}_\tau\|_2^2 ,
  \label{eq:drift-map-obj}
\end{equation}
where $\|\ve{A}\|_{\ve{M}}^2:=\Trace(\ve{A}^{\tr}\ve{M}\ve{A})$. The first term is the prior,
the second penalizes deviations from the deterministic drift
$\ve{S}_\tau\approx\ve{D}_\tau\ve{S}_{\tau-1}$, and the third is the observation loss at the
current key. Thus $\ve{\Omega}_\tau$ controls how much the memory is allowed to deviate from the
drifted previous memory: small $\ve{\Omega}_\tau$ enforces nearly deterministic transport, while
large $\ve{\Omega}_\tau$ allows abrupt changes.

The online filter is the one-step version of this dynamic least-squares problem. Given the
posterior at $t-1$, first predict
\begin{equation}
  \widehat{\ve{S}}_t=\ve{D}_t\ve{S}_{t-1},
  \qquad
  \widehat{\ve{P}}_t=\ve{D}_t\ve{P}_{t-1}\ve{D}_t^{\tr}+\ve{\Omega}_t .
  \label{eq:drift-filter-predict}
\end{equation}
Then the filtered memory is the minimizer
\begin{equation}
  \ve{S}_t
  =
  \argmin_{\ve{S}}
  \frac12\|\ve{S}-\widehat{\ve{S}}_t\|_{\widehat{\ve{P}}_t^{-1}}^2
  +\frac{1}{2r_t}
  \|\ve{S}^{\tr}\ve{k}_t-\ve{v}_t\|_2^2 .
  \label{eq:drift-filter-obj}
\end{equation}
Solving \eqref{eq:drift-filter-obj} gives the Kalman residual write in
\eqref{eq:mem-kf}. Compared with static RLS, the normal equations no longer accumulate around a
single $\ve{S}$; the batch objective is block-tridiagonal in the trajectory
$\ve{S}_{0:t}$, and the online solution carries the predicted covariance
$\widehat{\ve{P}}_t$ forward through the drift.

% =============================================================================
\subsection{Exact Kalman Delta Network}
\label{app:exact-kla}

Both diagonal KDN (\S\ref{sec:kla}) and isotropic KDN (\S\ref{app:iso-kla}) restrict the
uncertainty state to make it scan-computable. Before either approximation, it is useful to
write down the \emph{exact} Kalman filter as a layer, both as the target the approximations aim
at and as a ground-truth reference for testing. This is the Kalman optimal update of
Proposition~\ref{prop:kalman-optimal-update} instantiated with the same learned parameters as
the other KDN variants, and it is what we implement first as a naive recurrence.

The only structural commitment is on the noise model: the memory transition is diagonal,
$\ve{D}_t=\diag(\ve{\alpha}_t)$, and the process noise is an \emph{additive diagonal} matrix
$\ve{\Omega}_t=\diag(\ve{\omega}_t)$ rather than a dense one, so each key channel accrues its own
independently parameterized drift. Crucially, the covariance $\ve{P}_t$ itself is kept
\emph{dense}: the rank-one measurement update makes $\ve{P}_t$ dense after the first token, and
we do not project it back to a diagonal. The gain therefore stays the exact anisotropic Kalman
gain \eqref{eq:kalman-gain}, not its diagonal surrogate \eqref{eq:diag-kalman-gain}. The cost is
that $\ve{P}_t\in\R^{d_k\times d_k}$ is carried and updated in full ($O(d_k^2)$ state and work
per token) and the gain depends on the accumulated posterior covariance, so the recurrence does
not collapse to a parallel scan --- exactly the obstruction analyzed in
Appendix~\ref{app:deltanet-wy-exact-kalman}. This is acceptable for a reference implementation
and for the diagonal/isotropic equivalence checks; it is not the trained fast path.

\paragraph{Prediction.}
With $\ve{D}_t=\diag(\ve{\alpha}_t)$ and additive diagonal process noise
$\ve{\Omega}_t=\diag(\ve{\omega}_t)$, the covariance prediction of \eqref{eq:kalman-predict}
stays exact and is a cheap entrywise map,
\begin{equation}
  \widehat{\ve{P}}_t=\ve{D}_t\ve{P}_{t-1}\ve{D}_t^{\tr}+\ve{\Omega}_t,
  \qquad
  \widehat{P}_{t,ij}
  =
  \alpha_{t,i}\alpha_{t,j}\,P_{t-1,ij}+\omega_{t,i}\,\mathbf{1}[i=j],
  \label{eq:exact-kla-predict}
\end{equation}
alongside the memory prediction $\widehat{\ve{S}}_t=\ve{D}_t\ve{S}_{t-1}$. Only the diagonal of
$\widehat{\ve{P}}_t$ receives process noise; the off-diagonal correlations are transported by
the transition but not injected with drift.

\paragraph{Gain, update, and read.}
The gain, residual write, and covariance downdate are the unmodified Kalman update
\eqref{eq:kalman-gain}--\eqref{eq:kalman-update}, now carrying the dense
$\widehat{\ve{P}}_t$:
\begin{equation}
  \ve{\kappa}_t
  =
  \frac{\widehat{\ve{P}}_t\ve{k}_t}{r_t+\ve{k}_t^{\tr}\widehat{\ve{P}}_t\ve{k}_t},
  \qquad
  \ve{S}_t
  =
  \widehat{\ve{S}}_t+\ve{\kappa}_t\big(\ve{v}_t-\widehat{\ve{S}}_t^{\tr}\ve{k}_t\big)^{\tr},
  \qquad
  \ve{P}_t
  =
  (\ve{I}-\ve{\kappa}_t\ve{k}_t^{\tr})\widehat{\ve{P}}_t,
  \label{eq:exact-kla-update}
\end{equation}
and the layer reads $\ve{o}_t=\ve{S}_t^{\tr}\ve{q}_t$. Unlike the fixed-gain delta rule, the
gain direction $\ve{\kappa}_t$ is not proportional to $\ve{k}_t$: the dense
$\widehat{\ve{P}}_t$ rotates and rescales the write so that well-measured directions are edited
less than uncertain ones.

\refstepcounter{definition}\label{prop:exact-kalman-linear-attention}
\begin{tcolorbox}[
  title={Proposition~\thedefinition: Exact Kalman Delta Network},
  colback=blue!7,
  colbacktitle=blue!22,
  colframe=blue!45!black,
  coltitle=black,
  fonttitle=\bfseries,
  boxrule=0.4pt,
  arc=1pt,
  left=5pt,
  right=5pt,
  top=5pt,
  bottom=5pt
]
The exact Kalman Delta Network is the Kalman optimal update of
Proposition~\ref{prop:kalman-optimal-update} run as a layer with a diagonal transition and an
additive diagonal process noise, keeping the covariance $\ve{P}_t$ dense. Given token
representation $\ve{x}_t$, key--value pair $(\ve{k}_t,\ve{v}_t)$, query $\ve{q}_t$, previous
memory $\ve{S}_{t-1}$, and previous covariance $\ve{P}_{t-1}$, the layer parameterizes a
diagonal memory transition, an additive diagonal process noise, and a token-dependent
observation noise,
\begin{equation}
\begin{aligned}
  \ve{\alpha}_t &= \sigma(\ve{W}_{\alpha}\ve{x}_t+\ve{b}_{\alpha}),
  &
  \ve{\omega}_t &= \omega_{\min}+\operatorname{softplus}(\ve{W}_{\omega}\ve{x}_t+\ve{b}_{\omega}),
  &
  r_t &= r_{\min}+\operatorname{softplus}(\ve{w}_{r}^{\tr}\ve{x}_t+b_r),
  \\[2pt]
  \ve{D}_t &= \diag(\ve{\alpha}_t),
  &
  \ve{\Omega}_t &= \diag(\ve{\omega}_t),
\end{aligned}
\label{eq:exact-kla-box-parameterization}
\end{equation}
initialized with $\ve{S}_0=\ve{0}$ and isotropic prior covariance $\ve{P}_0=\mu^{-1}\ve{I}$. It
then runs the exact recurrence
\begin{equation}
\begin{aligned}
  \text{predict:}\quad
  \widehat{\ve{S}}_t &= \ve{D}_t\ve{S}_{t-1},
  &
  \widehat{\ve{P}}_t &= \ve{D}_t\ve{P}_{t-1}\ve{D}_t^{\tr}+\ve{\Omega}_t,
  \\[3pt]
  \text{gain:}\quad
  \ve{\kappa}_t &= \frac{\widehat{\ve{P}}_t\ve{k}_t}{r_t+\ve{k}_t^{\tr}\widehat{\ve{P}}_t\ve{k}_t},
  \\[3pt]
  \text{update:}\quad
  \ve{S}_t &= \widehat{\ve{S}}_t+\ve{\kappa}_t\big(\ve{v}_t-\widehat{\ve{S}}_t^{\tr}\ve{k}_t\big)^{\tr},
  &
  \ve{P}_t &= (\ve{I}-\ve{\kappa}_t\ve{k}_t^{\tr})\widehat{\ve{P}}_t,
  \\[3pt]
  \text{read:}\quad
  \ve{o}_t &= \ve{S}_t^{\tr}\ve{q}_t .
\end{aligned}
\label{eq:exact-kla-box-recurrence}
\end{equation}
Because $\ve{P}_t$ is dense and the gain depends on it, the recurrence is $O(d_k^2)$ per token
and is not a parallel scan; it is the ground-truth filter that diagonal KDN (\S\ref{sec:kla})
and isotropic KDN (\S\ref{app:iso-kla}) approximate, and the naive recurrent version we
implement first.
\end{tcolorbox}

\paragraph{Implementation notes.}
For a numerically robust reference, use the symmetric (Joseph) form of the covariance downdate,
which stays symmetric and positive semidefinite for any $\ve{\kappa}_t$:
\begin{equation}
  \ve{P}_t
  =
  (\ve{I}-\ve{\kappa}_t\ve{k}_t^{\tr})\widehat{\ve{P}}_t(\ve{I}-\ve{\kappa}_t\ve{k}_t^{\tr})^{\tr}
  +r_t\,\ve{\kappa}_t\ve{\kappa}_t^{\tr},
  \label{eq:exact-kla-joseph}
\end{equation}
which equals the posterior covariance $\ve{P}_t(\ve{\kappa}_t)$ minimized in
\eqref{eq:kalman-tracking-objective}; optionally re-symmetrize
$\ve{P}_t\leftarrow\tfrac12(\ve{P}_t+\ve{P}_t^{\tr})$. Two equivalence checks fall out directly:
projecting $\ve{P}_t$ to its diagonal $\diag(\ve{p}_t)$ after each step recovers the diagonal
predicted variance $\widehat{\ve{p}}_t=\ve{\alpha}_t^2\odot\ve{p}_{t-1}+\ve{\omega}_t$ of
\S\ref{sec:kla}, and averaging that diagonal recovers the isotropic scalar $\widehat{b}_t$ of
\S\ref{app:iso-kla}. To align the dense reference with KDN's standardized-coordinate Fisher
increment \eqref{eq:app-vector-update} for $\ell_2$-normalized keys, use the calibrated
measurement precision $d_k/r_t$ (equivalently standardized keys
$\widetilde{\ve{k}}_t=\sqrt{d_k}\,\ve{k}_t$) in the gain.

% =============================================================================
\subsection{Isotropic Kalman Delta Network}
\label{app:iso-kla}

Consider the most restrictive covariance family, the isotropic family
\begin{equation}
  \ve{P}_t\in\{b\ve{I}: b\in\R_+\},
  \qquad
  \ve{P}_t=b_t\ve{I},
  \qquad
  \ve{C}_t=\ve{P}_t^{-1}=c_t\ve{I},
  \quad c_t=b_t^{-1},
  \label{eq:iso-cov-family}
\end{equation}
which keeps a single scalar uncertainty $b_t$ (equivalently a scalar information $c_t=b_t^{-1}$)
for the entire key space. We derive the resulting update by specializing the diagonal KDN
recursion to this family with an independently parameterized additive process noise: predict
the scalar uncertainty through the transition, apply the scalar measurement update, and read
off the induced gain. The price of using an independent $\omega_t$ is that the information recursion
is no longer affine in $c_t$; it is a M\"obius map, which is still associative through a
$2\times2$ matrix scan.

\paragraph{Prediction.}
With a diagonal transition $\ve{D}_t=\diag(\ve{\alpha}_t)$ and additive process covariance
$\omega_t\ve{I}$, the covariance prediction is
$\widehat{\ve{P}}_t=b_{t-1}\ve{D}_t^2+\omega_t\ve{I}$. Projecting this diagonal matrix back to the
isotropic family by matching average variance gives
\begin{equation}
  \widehat{b}_t
  =
  a_t b_{t-1}+\omega_t,
  \qquad
  a_t=\frac{1}{d_k}\sum_{i=1}^{d_k}\ve{\alpha}_{t,i}^2,
  \qquad
  \widehat{c}_t=\frac{c_{t-1}}{a_t+\omega_t c_{t-1}}.
  \label{eq:iso-kla-predict}
\end{equation}
If $\ve{\alpha}_t$ is uniform across channels, this projection is exact. If it is channel-wise,
the memory transition can still be KDA's $\ve{D}_t=\diag(\ve{\alpha}_t)$, but the uncertainty
used by the scalar gain sees only the average contraction $a_t$. A diagonal process covariance
could be used instead by replacing $\omega_t$ with its average diagonal value.

\paragraph{Information update and associative scan.}
First consider the dimension-stable diagonal information update. Because an $\ell_2$-normalized
key spreads unit energy over $d_k$ coordinates, the scalar feature for channel $i$ is the
standardized coordinate $\widetilde{k}_{t,i}=\sqrt{d_k}\,\ve{k}_{t,i}$. The corresponding
per-channel Fisher increment is therefore
\begin{equation}
  \ve{c}_{t,i}
  =
  \widehat{\ve{c}}_{t,i}+d_k\frac{\ve{k}_{t,i}^2}{r_t},
  \qquad i=1,\dots,d_k .
  \label{eq:iso-kla-diag-info-update}
\end{equation}
For spread-out normalized keys this is $O(1/r_t)$ per channel in expectation, instead of
shrinking as $1/(d_k r_t)$. The isotropic tracker is the scalar average of these calibrated
diagonal precisions. Since the predicted covariance is already isotropic, its scalar information
is $\widehat{c}_t$; averaging the measurement increment and using
$\sum_i\ve{k}_{t,i}^2=\|\ve{k}_t\|_2^2$ gives
\begin{equation}
  c_t
  =
  \widehat{c}_t+\frac{1}{d_k}\sum_{i=1}^{d_k}d_k\frac{\ve{k}_{t,i}^2}{r_t}
  =
  \widehat{c}_t+\frac{\|\ve{k}_t\|_2^2}{r_t},
  \qquad
  c_0=\mu,
  \label{eq:iso-kla-info-update}
\end{equation}
which is the scalar measurement update. Thus the IsoKLA update is not an additional scaling
rule; it is the isotropic average of the same standardized-coordinate Fisher increments used by
the diagonal tracker. Equivalently, the isotropic projection of the calibrated dense Fisher
matrix $(d_k/r_t)\ve{k}_t\ve{k}_t^{\tr}$ is its trace average:
\begin{equation}
  u_t
  =
  \frac{1}{d_k}\Trace\!\left(\frac{d_k}{r_t}\ve{k}_t\ve{k}_t^{\tr}\right)
  =
  \frac{\|\ve{k}_t\|_2^2}{r_t}.
  \label{eq:iso-kla-trace-projection}
\end{equation}
This is different from summing all entries of $\ve{k}_t\ve{k}_t^{\tr}$, which gives
$(\ve{1}^{\tr}\ve{k}_t)^2$. That all-entry sum depends on coordinate signs and the chosen basis:
a key can have large norm but nearly zero coordinate sum, so it would incorrectly assign almost
no information to the measurement. Keeping the full rank-one matrix is the dense Kalman update;
once we restrict to an isotropic scalar state, the trace average is the invariant scalar
projection.

Combining the prediction and update gives a linear-fractional recurrence in the previous
information state. With
\begin{equation}
  u_t=\frac{\|\ve{k}_t\|_2^2}{r_t},
  \label{eq:iso-kla-u}
\end{equation}
we have
\begin{equation}
  c_t
  =
  \frac{c_{t-1}}{a_t+\omega_t c_{t-1}}+u_t
  =
  \frac{(1+u_t\omega_t)c_{t-1}+u_t a_t}{\omega_t c_{t-1}+a_t}.
  \label{eq:iso-kla-mobius}
\end{equation}
Thus each token defines a M\"obius map
\begin{equation}
  c_t=f_t(c_{t-1}),
  \qquad
  f_t(c)=\frac{A_t c+B_t}{C_t c+D_t},
  \qquad
  \ve{M}_t=
  \begin{bmatrix}
    A_t & B_t\\
    C_t & D_t
  \end{bmatrix}
  =
  \begin{bmatrix}
    1+u_t\omega_t & u_t a_t\\
    \omega_t & a_t
  \end{bmatrix}.
  \label{eq:iso-kla-mobius-matrix}
\end{equation}
M\"obius maps compose by matrix multiplication:
\begin{equation}
  f_t\circ f_{t-1}\circ\cdots\circ f_1
  \quad\Longleftrightarrow\quad
  \ve{M}_t\ve{M}_{t-1}\cdots\ve{M}_1 .
  \label{eq:iso-kla-mobius-compose}
\end{equation}
Equivalently, in homogeneous coordinates $c_t=n_t/d_t$ and
$[n_t,d_t]^{\tr}=\ve{M}_t[n_{t-1},d_{t-1}]^{\tr}$. Therefore all prefix information states can
be computed by an associative scan over $2\times2$ matrices, with $c_0=\mu$.

\paragraph{Gain.}
The predicted uncertainty is now isotropic, $\widehat{\ve{P}}_t=\widehat{b}_t\ve{I}$ with
$\widehat{b}_t=1/\widehat{c}_t$ from \eqref{eq:iso-kla-predict}. Substituting into the Kalman gain
$\ve{\kappa}_t=\widehat{\ve{P}}_t\ve{k}_t/(r_t+\ve{k}_t^{\tr}\widehat{\ve{P}}_t\ve{k}_t)$, the
covariance acts as a scalar and the gain collapses to the fixed delta-rule direction:
\begin{equation}
  \ve{\kappa}_t
  =
  \frac{\widehat{b}_t}{r_t+\widehat{b}_t\|\ve{k}_t\|_2^2}\,\ve{k}_t
  \equiv
  \beta_t\ve{k}_t .
  \label{eq:iso-kla-gain}
\end{equation}

\paragraph{Memory update.}
With $\beta_t$ from \eqref{eq:iso-kla-gain}, the memory update is exactly the
fixed-gain residual form
\begin{equation}
  \ve{S}_t
  =
  (\ve{I}-\beta_t\ve{k}_t\ve{k}_t^{\tr})\ve{D}_t\ve{S}_{t-1}
  +\beta_t\ve{k}_t\ve{v}_t^{\tr}.
  \label{eq:iso-kla-memory}
\end{equation}
Thus isotropic KDN can adapt the scalar write strength $\beta_t$ through the Kalman uncertainty
scan, but --- keeping only one scalar for the entire key space --- it cannot assign different
write strengths to different key channels; an anisotropic gain is recovered only by relaxing to
the diagonal family.

\refstepcounter{definition}\label{prop:isotropic-kalman-linear-attention}
\begin{tcolorbox}[
  title={Proposition~\thedefinition: Isotropic Kalman Delta Network},
  colback=blue!7,
  colbacktitle=blue!22,
  colframe=blue!45!black,
  coltitle=black,
  fonttitle=\bfseries,
  boxrule=0.4pt,
  arc=1pt,
  left=5pt,
  right=5pt,
  top=5pt,
  bottom=5pt
]
The isotropic Kalman Delta Network is the scalar-uncertainty specialization of KDN. Given
token representation $\ve{x}_t$, key--value pair $(\ve{k}_t,\ve{v}_t)$, and previous memory
$\ve{S}_{t-1}$, the layer parameterizes a diagonal memory transition, an additive process noise,
and a token-dependent observation noise,
\begin{equation}
\begin{aligned}
  \ve{\alpha}_t &= \sigma(\ve{W}_{\alpha}\ve{x}_t+\ve{b}_{\alpha}),
  &
  \omega_t &= \omega_{\min}+\operatorname{softplus}(\ve{w}_{\omega}^{\tr}\ve{x}_t+b_\omega),
  &
  r_t &= r_{\min}+\operatorname{softplus}(\ve{w}_{r}^{\tr}\ve{x}_t+b_r),
  \\[2pt]
  a_t &= \frac{1}{d_k}\sum_i\ve{\alpha}_{t,i}^2,
  &
  \ve{D}_t &= \diag(\ve{\alpha}_t).
\end{aligned}
\label{eq:iso-kla-box-parameterization}
\end{equation}
It then computes scalar prefix information statistics by a M\"obius scan:
\begin{equation}
  u_t=\frac{\|\ve{k}_t\|_2^2}{r_t},
  \qquad
  \ve{M}_t=
  \begin{bmatrix}
    1+u_t\omega_t & u_t a_t\\
    \omega_t & a_t
  \end{bmatrix},
  \qquad
  \begin{bmatrix}n_t\\ d_t\end{bmatrix}
  =
  \ve{M}_t
  \begin{bmatrix}n_{t-1}\\ d_{t-1}\end{bmatrix},
  \qquad
  c_t=\frac{n_t}{d_t},
  \qquad
  \begin{bmatrix}n_0\\ d_0\end{bmatrix}
  =
  \begin{bmatrix}\mu\\ 1\end{bmatrix}.
  \label{eq:iso-kla-box-mobius}
\end{equation}
The gain at token $t$ uses the exclusive prefix state $c_{t-1}$:
\begin{equation}
\begin{aligned}
  \widehat{c}_t
  &= \frac{c_{t-1}}{a_t+\omega_t c_{t-1}},
  &
  \widehat{b}_t
  &= \frac{1}{\widehat{c}_t},
  &
  \beta_t
  &= \frac{\widehat{b}_t}{r_t+\widehat{b}_t\|\ve{k}_t\|_2^2}.
\end{aligned}
\label{eq:iso-kla-box-gain}
\end{equation}
With $\beta_t$ available before the memory scan, the memory update is the fixed-direction
affine scan
\begin{equation}
  \ve{S}_t
  =
  (\ve{I}-\beta_t\ve{k}_t\ve{k}_t^{\tr})\ve{D}_t\ve{S}_{t-1}
  +\beta_t\ve{k}_t\ve{v}_t^{\tr}.
  \label{eq:iso-kla-box-memory}
\end{equation}
Thus ISO-KDN preserves the Kalman idea of adapting write strength from uncertainty, but the
isotropic covariance restriction forces the gain direction to be $\ve{k}_t$.
\end{tcolorbox}

Therefore ISO-KDN sits between fixed-gain delta rules and diagonal KDN: it keeps a
Kalman-style scalar uncertainty and token-dependent observation noise, but because
$\ve{P}_t$ is isotropic, it always writes along $\ve{k}_t$ with one scalar strength. Diagonal
Diagonal KDN strictly generalizes this by replacing the scalar $b_t$ with a vector
$\ve{p}_t$, allowing the gain to be anisotropic.

% =============================================================================
\subsection{Diagonal information: from the covariance recursion to the M\"obius scan}
\label{app:diag-info}

This section derives, step by step, the predicted uncertainty $\widehat{\ve{p}}_t$ of
\S\ref{sec:kla} from the diagonal covariance recursion and identifies the only approximation:
the dense measurement information is projected back to a diagonal state. With an independently
parameterized process noise, the resulting information update is a M\"obius map; this remains
parallelizable because M\"obius maps compose by $2\times2$ matrix multiplication. All products,
inverses, and ratios are coordinatewise after the diagonal projection; we then write a single
channel and drop the index.

\paragraph{Diagonal state.}
Restrict the key-space covariance of the associative-memory filter \eqref{eq:mem-kf} to
diagonal, $\ve{P}_t=\diag(\ve{p}_t)$, with transition $\ve{D}_t=\diag(\ve{\alpha}_t)$ and process
noise $\ve{\Omega}_t=\diag(\ve{\omega}_t)$. Per channel, track the variance $p_t$ and its reciprocal, the
information $c_t=1/p_t$.

\paragraph{Step 1 --- predict (the master recursion).}
The covariance predict $\widehat{\ve{P}}_t=\ve{D}_t\ve{P}_{t-1}\ve{D}_t^{\tr}+\ve{\Omega}_t$ preserves
diagonality, and per channel is
\begin{equation}
  \boxed{\;\widehat{p}_t=\alpha_t^2\,p_{t-1}+\omega_t\;}
  \label{eq:app-predict}
\end{equation}
(the diagonal form of \eqref{eq:diag-covariance-prediction}). This is the \emph{only} exact
predicted variance; every expression below is \eqref{eq:app-predict} written in different
variables.

\paragraph{Step 2 --- update, in information form.}
Restore vector notation for the measurement step. Given the predicted diagonal information
$\widehat{\ve{C}}_t=\diag(\widehat{\ve{c}}_t)$, the exact linear-Gaussian information update is
\begin{equation}
  \ve{C}_t^{\mathrm{exact}}
  =
  \widehat{\ve{C}}_t+\frac{1}{r_t}\ve{k}_t\ve{k}_t^{\tr}.
  \label{eq:app-exact-info-update}
\end{equation}
The added Fisher information is rank one and generally dense, so keeping only a diagonal vector
$\ve{c}_t$ is a projection approximation. The literal diagonal projection would add
$\ve{k}_t\odot\ve{k}_t/r_t$. For $\ell_2$-normalized keys, however, a raw coordinate has typical
squared size $1/d_k$, so that update gives each channel only $O(1/(d_k r_t))$ information.
Using the standardized coordinate $\widetilde{k}_{t,i}=\sqrt{d_k}\,\ve{k}_{t,i}$ gives the
calibrated diagonal update
\begin{equation}
  \ve{c}_t=\widehat{\ve{c}}_t+\frac{d_k}{r_t}\,\ve{k}_t\odot\ve{k}_t.
  \label{eq:app-vector-update}
\end{equation}
Equivalently, per channel,
\begin{equation}
  c_t=\widehat{c}_t+d_k\frac{k_t^2}{r_t},
  \qquad
  \widehat{c}_t:=\frac{1}{\widehat{p}_t}.
  \label{eq:app-update}
\end{equation}
For normalized keys with roughly uniform coordinate energy, this gives
$\mathbb{E}[d_k k_t^2/r_t]\approx 1/r_t$ per channel. The predicted information
$\widehat{c}_t$ is still \emph{defined} as $1/\widehat{p}_t$; it is not a separate quantity.

\paragraph{Step 3 --- the recursion.}
Insert \eqref{eq:app-predict} into $\widehat{c}_t=1/\widehat{p}_t$ and use $p_{t-1}=1/c_{t-1}$:
\begin{equation}
  \widehat{c}_t
  =\frac{1}{\alpha_t^2\,p_{t-1}+\omega_t}
  =\frac{1}{\alpha_t^2/c_{t-1}+\omega_t}
  =\frac{c_{t-1}}{\alpha_t^2+\omega_t\,c_{t-1}},
  \label{eq:app-chat}
\end{equation}
so the full posterior recursion is
\begin{equation}
  c_t=\frac{c_{t-1}}{\alpha_t^2+\omega_t\,c_{t-1}}+d_k\frac{k_t^2}{r_t}.
  \label{eq:app-cfull}
\end{equation}
For a general $\omega_t$ this is a \emph{M\"obius} (linear-fractional) map of $c_{t-1}$,
$c_t=(A_tc_{t-1}+B_t)/(C_tc_{t-1}+D_t)$, with
\begin{equation}
  \begin{bmatrix}A_t&B_t\\[2pt] C_t&D_t\end{bmatrix}
  =\begin{bmatrix}1+d_k\tfrac{k_t^2}{r_t}\,\omega_t & d_k\tfrac{k_t^2}{r_t}\,\alpha_t^2\\[3pt] \omega_t & \alpha_t^2\end{bmatrix}.
  \label{eq:app-mobius}
\end{equation}
M\"obius maps compose by matrix multiplication, so $c_T$ is the action of the product
$\ve{M}_T\cdots\ve{M}_1$ on $c_0$. When $\omega_t$ is input-only the matrices are input-only, the
product is associative, and it parallelizes as the scan in \eqref{eq:diag-info-scan} (at higher
constant cost than an affine scan, and needing per-chunk renormalization since the entries can
span a wide range). This is the diagonal-information recursion used by KDN with freely
parameterized additive process noise.

\paragraph{Step 4 --- the special choice that linearizes the scan.}
Equation~\eqref{eq:app-cfull} is affine in $c_{t-1}$ iff the product $\omega_t\,c_{t-1}$ is
state-independent, i.e.\ iff $\omega_t\propto p_{t-1}$. Writing this as
\begin{equation}
  \omega_t=(\lambda_t^{-1}-\alpha_t^2)\,p_{t-1}
  \qquad\Longleftrightarrow\qquad
  \lambda_t=\frac{1}{\alpha_t^2+\omega_t/p_{t-1}}
  \label{eq:app-lambda}
\end{equation}
(the identity \eqref{eq:app-lambda}) makes $\alpha_t^2+\omega_t\,c_{t-1}=\lambda_t^{-1}$, so
\eqref{eq:app-chat}--\eqref{eq:app-cfull} collapse to the affine scan
\begin{equation}
  \widehat{c}_t=\lambda_t\,c_{t-1},
  \qquad
  c_t=\lambda_t\,c_{t-1}+d_k\frac{k_t^2}{r_t},
  \label{eq:app-affine}
\end{equation}
which is the cheaper affine-scan restriction. Substituting the \emph{same} $\omega_t$ back into the master
recursion \eqref{eq:app-predict},
\begin{equation}
  \widehat{p}_t=\alpha_t^2\,p_{t-1}+\omega_t=\lambda_t^{-1}\,p_{t-1}=\frac{1}{\widehat{c}_t},
  \label{eq:app-phat}
\end{equation}
i.e.\ \eqref{eq:prefix-diag-uncertainty}. The transition term $\alpha_t^2$ is not dropped; it is
exactly cancelled by this $\omega_t$. The small prior $\mu$ enters only through the initialization
$c_0=\mu$, which keeps the information state positive at the start of the scan.

\paragraph{Recovering KDA.}
KDN reduces to KDA when the transition remains KDA's diagonal transition and the predicted
uncertainty becomes isotropic. In the diagonal information form, this means
\begin{equation}
  \ve{D}_t=\diag(\ve{\alpha}_t),
  \qquad
  \widehat{\ve{p}}_t=b_t\ve{1}
  \quad\Longleftrightarrow\quad
  \widehat{\ve{c}}_t=b_t^{-1}\ve{1}
  \label{eq:app-kda-isotropic-condition}
\end{equation}
for some scalar $b_t>0$. Since KDN parameterizes an additive process noise and uses
$\widehat{p}_{t,i}=\alpha_{t,i}^2p_{t-1,i}+\omega_{t,i}$, the isotropic condition is
\begin{equation}
  \alpha_{t,i}^2p_{t-1,i}+\omega_{t,i}
  =
  b_t,
  \qquad\text{equivalently}\qquad
  \omega_{t,i}=b_t-\alpha_{t,i}^2p_{t-1,i},
  \label{eq:app-kda-omega-condition}
\end{equation}
for every channel $i$, with the validity condition $\omega_{t,i}\ge 0$. Under
\eqref{eq:app-kda-isotropic-condition}, the KDN gain collapses to
\begin{equation}
  \ve{\kappa}_t
  =
  \frac{b_t}{r_t+b_t\|\ve{k}_t\|_2^2}\,\ve{k}_t
  \equiv \beta_t\ve{k}_t,
  \label{eq:app-kda-gain-collapse}
\end{equation}
and the memory update becomes the vanilla KDA residual update
\begin{equation}
  \ve{S}_t
  =
  (\ve{I}-\beta_t\ve{k}_t\ve{k}_t^{\tr})\diag(\ve{\alpha}_t)\ve{S}_{t-1}
  +\beta_t\ve{k}_t\ve{v}_t^{\tr}.
  \label{eq:app-kda-recovered}
\end{equation}
For normalized keys, $\beta_t=b_t/(r_t+b_t)$.

\paragraph{Consequences.}
\begin{itemize}[leftmargin=1.4em,itemsep=3pt,topsep=2pt]
  \item Equation~\eqref{eq:app-predict}, $\widehat{p}_t=\alpha_t^2 p_{t-1}+\omega_t$, is the master
  recursion; \eqref{eq:app-chat} is the same recursion written in information variables.
  \item $\widehat{p}_t=1/\widehat{c}_t$ is not an additional modeling term. Since
  $\widehat{c}_t=c_{t-1}/(\alpha_t^2+\omega_t c_{t-1})$, its inverse is exactly
  $\alpha_t^2p_{t-1}+\omega_t$. Adding $\omega_t$ again would double-count the process noise.
  \item Free, input-only, $p$-independent $\omega_t$ gives the M\"obius scan
  \eqref{eq:app-mobius}; the older $\lambda$-form $\omega_t\propto p_{t-1}$ is the special case that
  collapses this to the affine scan \eqref{eq:app-affine}.
\end{itemize}


% =============================================================================
\section{Can the DeltaNet WY trick parallelize exact Kalman?}
\label{app:deltanet-wy-exact-kalman}

The chunkwise DeltaNet algorithm of \citet{yang2024parallelizing} is a useful comparison point,
but it does not directly remove the covariance recursion of the exact Kalman filter. The reason is
where the associativity lives.

\paragraph{What the DeltaNet trick parallelizes.}
Up to transpose conventions, the fixed-gain delta update has the form
\begin{equation}
  \ve{S}_t
  =
  (\ve{I}-\beta_t\ve{k}_t\ve{k}_t^{\tr})\ve{S}_{t-1}
  +\beta_t\ve{k}_t\ve{v}_t^{\tr}.
  \label{eq:app-deltanet-fixed-gain}
\end{equation}
The erase direction $\beta_t\ve{k}_t$ is known from the current token before the memory state is
updated. Therefore the within-chunk transition matrices
$\ve{I}-\beta_t\ve{k}_t\ve{k}_t^{\tr}$ are input-known rank-one maps. DeltaNet exploits this by
representing their product with a compact WY/UT form.

To see the algebra, use DeltaNet's row-state convention within a chunk of length $C$:
\begin{equation}
  \ve{S}_r
  =
  \ve{S}_{r-1}(\ve{I}-\beta_r\ve{k}_r\ve{k}_r^{\tr})
  +\beta_r\ve{v}_r\ve{k}_r^{\tr},
  \qquad r=1,\ldots,C,
  \label{eq:app-deltanet-row-update}
\end{equation}
with chunk input state $\ve{S}_0$. Write
$\ve{H}_r=\ve{I}-\beta_r\ve{k}_r\ve{k}_r^{\tr}$. The unrolled state has the form
\begin{equation}
  \ve{S}_r=\ve{S}_0\ve{P}_r+\ve{R}_r,
  \qquad
  \ve{P}_r=\ve{H}_1\cdots\ve{H}_r.
  \label{eq:app-deltanet-unroll}
\end{equation}
The compact representation is that both $\ve{P}_r$ and the additive term $\ve{R}_r$ can be
written using only vectors:
\begin{equation}
  \ve{P}_r
  =
  \ve{I}-\sum_{i=1}^r \ve{w}_i\ve{k}_i^{\tr},
  \qquad
  \ve{R}_r
  =
  \sum_{i=1}^r \ve{u}_i\ve{k}_i^{\tr}.
  \label{eq:app-deltanet-compact-pr}
\end{equation}
Induction gives the recurrences
\begin{equation}
\begin{aligned}
  \ve{w}_r
  &=
  \beta_r
  \left(
    \ve{k}_r-\sum_{i<r}\ve{w}_i(\ve{k}_i^{\tr}\ve{k}_r)
  \right),
  \\
  \ve{u}_r
  &=
  \beta_r
  \left(
    \ve{v}_r-\sum_{i<r}\ve{u}_i(\ve{k}_i^{\tr}\ve{k}_r)
  \right).
\end{aligned}
\label{eq:app-deltanet-wu-rec}
\end{equation}
These equations still look sequential, but the dependency is only lower triangular within the
chunk. Stack the row vectors into
$\ve{K},\ve{W}\in\R^{C\times d_k}$ and
$\ve{V},\ve{U}\in\R^{C\times d_v}$, and let
$\ve{B}=\diag(\ve{\beta})$. If
\begin{equation}
  \ve{L}=\operatorname{tril}(\ve{B}\ve{K}\ve{K}^{\tr},-1),
  \label{eq:app-deltanet-lower}
\end{equation}
then \eqref{eq:app-deltanet-wu-rec} is exactly the pair of triangular systems
\begin{equation}
  (\ve{I}_C+\ve{L})\ve{W}=\ve{B}\ve{K},
  \qquad
  (\ve{I}_C+\ve{L})\ve{U}=\ve{B}\ve{V}.
  \label{eq:app-deltanet-tri-systems}
\end{equation}
Equivalently,
\begin{equation}
  \ve{T}=(\ve{I}_C+\ve{L})^{-1}\ve{B},
  \qquad
  \ve{W}=\ve{T}\ve{K},
  \qquad
  \ve{U}=\ve{T}\ve{V}.
  \label{eq:app-deltanet-ut}
\end{equation}
This is the UT form used by \citet{yang2024parallelizing}. The only solve is over the
$C\times C$ lower-triangular matrix $\ve{I}_C+\ve{L}$; after it is solved, the rest of the
chunk is standard matrix multiplication.

Substituting \eqref{eq:app-deltanet-compact-pr} into
\eqref{eq:app-deltanet-unroll} gives
\begin{equation}
  \ve{S}_r
  =
  \ve{S}_0+\sum_{i=1}^r
  \left(\ve{u}_i-\ve{S}_0\ve{w}_i\right)\ve{k}_i^{\tr}.
  \label{eq:app-deltanet-chunk-state}
\end{equation}
Thus, with $\widetilde{\ve{U}}=\ve{U}-\ve{W}\ve{S}_0^{\tr}$ and causal mask
$\ve{M}_C$, the whole chunk is
\begin{equation}
  \ve{S}_C=\ve{S}_0+\widetilde{\ve{U}}^{\tr}\ve{K},
  \qquad
  \ve{O}
  =
  \ve{Q}\ve{S}_0^{\tr}
  +(\ve{Q}\ve{K}^{\tr}\odot\ve{M}_C)\widetilde{\ve{U}}.
  \label{eq:app-deltanet-chunk-parallel}
\end{equation}
This is the parallelizable form: no per-token hidden matrices $\ve{S}_1,\ldots,\ve{S}_C$ are
materialized, and the recurrent-looking inner update has become one triangular solve plus
matrix multiplications.

The same idea applies to any fixed-gain residual memory update whose erase/write directions are
available before the memory scan. In particular, once KDN has computed a sequence of gains
$\ve{\kappa}_t$, the memory update
\begin{equation}
  \ve{S}_t
  =
  (\ve{I}-\ve{\kappa}_t\ve{k}_t^{\tr})\ve{D}_t\ve{S}_{t-1}
  +\ve{\kappa}_t\ve{v}_t^{\tr}
  \label{eq:app-kla-memory-fixed-kappa}
\end{equation}
is again an affine fast-weight recurrence with input-known coefficients. DeltaNet-style chunking
can then parallelize the \emph{memory} update.

\paragraph{Where exact Kalman differs.}
In the exact filter, however, the gain is not input-known:
\begin{equation}
  \ve{\kappa}_t
  =
  \frac{\widehat{\ve{P}}_t\ve{k}_t}
       {r_t+\ve{k}_t^{\tr}\widehat{\ve{P}}_t\ve{k}_t},
  \qquad
  \widehat{\ve{P}}_t
  =
  \ve{D}_t\ve{P}_{t-1}\ve{D}_t^{\tr}+\ve{\Omega}_t .
  \label{eq:app-exact-kalman-state-dependent-gain}
\end{equation}
The coefficient of the rank-one memory map in \eqref{eq:app-kla-memory-fixed-kappa} depends on
the accumulated posterior covariance. Thus the triangular WY system cannot be formed from
$(\ve{k}_t,\ve{v}_t,\ve{D}_t,\ve{\Omega}_t,r_t)$ alone: one must first solve the uncertainty
recursion that determines $\widehat{\ve{P}}_t$.

Equivalently, the exact covariance update is a Riccati recursion,
\begin{equation}
  \ve{P}_t
  =
  \widehat{\ve{P}}_t
  -
  \frac{\widehat{\ve{P}}_t\ve{k}_t\ve{k}_t^{\tr}\widehat{\ve{P}}_t}
       {r_t+\ve{k}_t^{\tr}\widehat{\ve{P}}_t\ve{k}_t}
  =
  \left(\widehat{\ve{P}}_t^{-1}+\frac{1}{r_t}\ve{k}_t\ve{k}_t^{\tr}\right)^{-1}.
  \label{eq:app-full-riccati}
\end{equation}
This is the part not handled by DeltaNet's WY representation. WY parallelizes products of
already-known rank-one memory transitions; it does not make the state-dependent Kalman gains
available for free.

\paragraph{Associative in principle, but not cheap.}
The full Riccati recursion still has algebraic structure. Assuming $\ve{D}_t$ is invertible, the
prediction step can be written as the matrix linear-fractional map
\begin{equation}
  \widehat{\ve{P}}_t
  =
  \big(\ve{D}_t\ve{P}_{t-1}
      +\ve{\Omega}_t\ve{D}_t^{-\tr}\big)
  \big(\ve{D}_t^{-\tr}\big)^{-1},
  \label{eq:app-predict-lft}
\end{equation}
and the measurement update as
\begin{equation}
  \ve{P}_t
  =
  \ve{P}_t^-
  \left(\frac{1}{r_t}\ve{k}_t\ve{k}_t^{\tr}\ve{P}_t^-+\ve{I}\right)^{-1}.
  \label{eq:app-measure-lft}
\end{equation}
Thus each token can be represented by a block matrix acting on covariance matrices through
\begin{equation}
  \ve{P}\mapsto(\ve{A}\ve{P}+\ve{B})(\ve{C}\ve{P}+\ve{D})^{-1},
  \qquad
  \begin{bmatrix}\ve{A}&\ve{B}\\ \ve{C}&\ve{D}\end{bmatrix}\in\R^{2d_k\times 2d_k}.
  \label{eq:app-matrix-mobius}
\end{equation}
These matrix M\"obius maps compose by block-matrix multiplication, exactly like the scalar
$2\times2$ maps in \eqref{eq:app-mobius}. Therefore an exact parallel scan exists in principle.
But its scan state is dense: composing two summaries costs cubic time in $d_k$ in the generic
case, and materializing prefix summaries costs quadratic memory per token. This is a very
different object from the DeltaNet chunk solve, whose within-chunk work is built from
$C\times C$ triangular systems and $d_k$-wide matrix multiplications.

\paragraph{Special cases and practical routes.}
The exact measurement update is simple in dense information form:
\begin{equation}
  \ve{C}_t
  =
  \widehat{\ve{C}}_t+\frac{1}{r_t}\ve{k}_t\ve{k}_t^{\tr}.
  \label{eq:app-dense-info-rank-one}
\end{equation}
The difficulty is the next prediction step when additive process noise is present. Process noise is
added in covariance space, so with $\ve{C}_{t-1}=\ve{P}_{t-1}^{-1}$ the predicted information is
\begin{equation}
  \widehat{\ve{C}}_t
  =
  \left(
    \ve{D}_t\ve{C}_{t-1}^{-1}\ve{D}_t^{\tr}+\ve{\Omega}_t
  \right)^{-1}.
  \label{eq:app-dense-info-predict}
\end{equation}
If $\ve{\Omega}_t=0$, this reduces to a linear information prediction,
\begin{equation}
  \ve{C}_t^-
  =
  \ve{D}_t^{-\tr}\ve{C}_{t-1}\ve{D}_t^{-1},
  \qquad
  \ve{C}_t
  =
  \ve{C}_t^-+\frac{1}{r_t}\ve{k}_t\ve{k}_t^{\tr}.
  \label{eq:app-no-drift-linear-info}
\end{equation}
For diagonal $\ve{D}_t=\diag(\ve{\alpha}_t)$, this is just entrywise scaling,
$\widehat{C}_{t,ij}=C_{t-1,ij}/(\alpha_{t,i}\alpha_{t,j})$, followed by the rank-one add
\eqref{eq:app-dense-info-rank-one}. Thus the no-drift dense information recursion is an affine
scan over dense matrices.

With additive drift $\ve{\Omega}_t\ne0$, the inverse in \eqref{eq:app-dense-info-predict} does
not distribute over the sum. Even if both $\ve{D}_t$ and $\ve{\Omega}_t$ are diagonal, a previous
rank-one measurement has already made $\ve{C}_{t-1}$ dense. Define
\begin{equation}
  \ve{G}_t=\ve{D}_t^{-\tr}\ve{C}_{t-1}\ve{D}_t^{-1}.
  \label{eq:app-dense-g}
\end{equation}
Woodbury gives the equivalent information prediction
\begin{equation}
  \widehat{\ve{C}}_t
  =
  \ve{G}_t
  -
  \ve{G}_t
  \left(\ve{\Omega}_t^{-1}+\ve{G}_t\right)^{-1}
  \ve{G}_t .
  \label{eq:app-dense-woodbury-predict}
\end{equation}
Although $\ve{\Omega}_t^{-1}$ is diagonal, $\ve{G}_t$ is dense, so
$\ve{\Omega}_t^{-1}+\ve{G}_t$ is dense and computing $\widehat{\ve{C}}_t$ requires a dense solve.
This is the precise sense in which additive diagonal process noise breaks the simple dense
information scan: the observation update remains a rank-one add, but the prediction step becomes
a dense matrix-fractional map.

\paragraph{Possible workarounds.}
There are several ways around this obstruction, but each changes either the exactness target or
the computational profile.
\begin{enumerate}[leftmargin=1.6em,itemsep=3pt,topsep=3pt]
  \item \textbf{Track dense covariance directly.}
  In covariance variables, additive diagonal process noise is easy:
  \begin{equation}
    \widehat{\ve{P}}_t=\ve{D}_t\ve{P}_{t-1}\ve{D}_t^{\tr}+\ve{\Omega}_t,
    \qquad
    \widehat{P}_{t,ij}
    =
    \alpha_{t,i}\alpha_{t,j}P_{t-1,ij}
    +\omega_{t,i}\mathbf{1}[i=j].
    \label{eq:app-dense-cov-predict}
  \end{equation}
  The measurement update is also an explicit rank-one downdate,
  \begin{equation}
    \ve{P}_t
    =
    \widehat{\ve{P}}_t
    -
    \frac{\widehat{\ve{P}}_t\ve{k}_t\ve{k}_t^{\tr}\widehat{\ve{P}}_t}
         {r_t+\ve{k}_t^{\tr}\widehat{\ve{P}}_t\ve{k}_t}.
    \label{eq:app-dense-cov-update}
  \end{equation}
  This gives an exact dense recurrent reference without matrix inversion, but the gain still
  depends on the updated covariance, so the sequence remains state-dependent and not a cheap
  DeltaNet-style chunk scan.

  \item \textbf{Remove additive drift.}
  Setting $\ve{\Omega}_t=0$ gives the affine dense-information scan
  \eqref{eq:app-no-drift-linear-info}. This is exact for a deterministic transition model, but it
  removes independent process drift. For contractive $\alpha_{t,i}<1$, the deterministic map
  shrinks covariance and therefore increases information unless additional modeling choices
  counteract it.

  \item \textbf{Use state-proportional process noise.}
  Choose process noise so that prediction is multiplicative in covariance,
  \begin{equation}
    \widehat{\ve{P}}_t
    =
    \lambda_t^{-1}\ve{D}_t\ve{P}_{t-1}\ve{D}_t^{\tr},
    \qquad
    \ve{\Omega}_t
    =
    (\lambda_t^{-1}-1)\ve{D}_t\ve{P}_{t-1}\ve{D}_t^{\tr}.
    \label{eq:app-state-proportional-drift}
  \end{equation}
  Then
  \begin{equation}
    \widehat{\ve{C}}_t
    =
    \lambda_t\ve{D}_t^{-\tr}\ve{C}_{t-1}\ve{D}_t^{-1},
    \label{eq:app-state-proportional-info}
  \end{equation}
  so the information scan is affine again. The price is semantic: $\ve{\Omega}_t$ is no longer an
  independently parameterized diagonal additive noise; it is tied to the current posterior
  covariance.

  \item \textbf{Use block-diagonal covariance.}
  Partition key channels into blocks and keep dense Kalman algebra only within each block:
  \begin{equation}
    \ve{C}_t^{(j)}
    =
    \widehat{\ve{C}}_t^{(j)}
    +
    \frac{1}{r_t}\ve{k}_t^{(j)}\ve{k}_t^{(j)\tr}.
    \label{eq:app-block-info-update}
  \end{equation}
  This preserves within-block correlations and replaces one full $d_k\times d_k$ solve by many
  small $b\times b$ solves or matrix M\"obius scans. It is approximate because the cross-block
  Fisher terms $\ve{k}_t^{(j)}\ve{k}_t^{(\ell)\tr}/r_t$ are dropped.

  \item \textbf{Use diagonal plus low rank information.}
  Maintain
  \begin{equation}
    \ve{C}_t=\diag(\ve{c}_t)+\ve{U}_t\ve{U}_t^{\tr}.
    \label{eq:app-diag-low-rank-info}
  \end{equation}
  The measurement update appends a rank-one direction. Woodbury can then keep predictions and
  gains tractable as long as the rank is controlled. This requires a truncation or compression
  rule, so it is a richer approximation rather than an exact fixed-size filter.
\end{enumerate}

The practical conclusion is simple: the WY trick and exact Kalman address different
bottlenecks. The WY trick is valuable after the gains are known; exact Kalman still needs a
tractable gain recursion. KDN chooses the diagonal restriction, where the Riccati map reduces to
independent scalar M\"obius scans. A natural middle ground is block-diagonal covariance: use
small $b\times b$ matrix M\"obius scans to compute blockwise gains, then feed those gains into a
chunked memory update. This keeps the exact Kalman algebra inside each block, while making the
cost scale with the chosen block size rather than the full head dimension.



\end{document}
